\documentclass[aspectratio=169,10pt]{beamer}
\usetheme{mekatronik}
\renewcommand{\coursetitle}{Mechatronics and Robotics Engineering (50 h)}
\setweek{1}
\setday{1}{Electrical Machines \& Actuators in Robots I}
\title[Machines \& Actuators in Robots I]{Electrical Machines and Actuators in Robots I}
\subtitle{Energy conversion and the actuator landscape \textperiodcentered{} DC, BLDC and PMSM motors\\ \textperiodcentered{} motor models and time constants \textperiodcentered{} stepper and induction motors}
\author{Mechatronics and Robotics Engineering Programme}
\institute{Week 1 \textperiodcentered{} Day 1}
\date{2 weeks \textperiodcentered{} 10 days \textperiodcentered{} 50 hours}

% small helpers used only in this deck
\newcommand{\figph}[2]{\includegraphics[#1]{../../figures/#2}}
\newcommand{\gofurther}[1]{{\scriptsize\color{mkGray}\textbf{Going further:} #1\par}}
\tikzset{
  mkblk/.style={draw=mkBlue,fill=mkPale,rounded corners=1.5pt,minimum height=7mm,minimum width=12mm,font=\scriptsize,align=center,line width=.6pt},
  mksum/.style={draw=mkBlue,circle,fill=white,minimum size=4.5mm,inner sep=0pt,line width=.6pt},
  mksig/.style={-{Stealth[length=2mm]},mkNavy,line width=.6pt},
  mklab/.style={font=\scriptsize,text=mkNavy},
}

\begin{document}

% =====================================================================
\begin{frame}[plain]\titlepage\end{frame}

% =====================================================================
\begin{frame}{Agenda and learning outcomes}
\begin{columns}[T]
\begin{column}{.52\textwidth}
\small
\renewcommand{\arraystretch}{1.15}
\begin{tabular}{@{}l p{4.7cm} r@{}}
\toprule
\textbf{Code} & \textbf{Session} & \textbf{Time}\\
\midrule
1.1 & Electromechanical energy conversion \& the actuator landscape in robots & 1.0\,h\\
1.2 & DC, BLDC and PMSM motors: structure and commutation & 1.5\,h\\
1.3 & Electrical and mechanical motor models & 1.5\,h\\
1.4 & Stepper and induction motors in automation & 1.0\,h\\
\midrule
 & \textbf{Net teaching time} & \textbf{5.0\,h}\\
\bottomrule
\end{tabular}\\[1mm]
{\scriptsize\color{mkGray} Breaks ($\approx$10 min) roughly every hour are additional.}
\end{column}
\begin{column}{.45\textwidth}
\begin{defbox}[By the end of today you can \dots]
\small
\begin{itemize}\setlength\itemsep{1pt}
\item \textbf{explain} how a motor turns current into torque
\item \textbf{compare} electric, pneumatic and hydraulic actuators for a robot
\item \textbf{sketch} how DC, BLDC and PMSM motors commutate
\item \textbf{write} the DC-motor equations and \textbf{calculate} its time constants
\item \textbf{choose} between stepper, induction and servo motors
\end{itemize}
\end{defbox}
\end{column}
\end{columns}
\end{frame}

% =====================================================================
\section{1.1 Energy conversion \& the actuator landscape}
% =====================================================================

\begin{frame}{A robot joint is a chain of energy converters}
\begin{columns}[T]
\begin{column}{.68\textwidth}
\centering
\begin{tikzpicture}[node distance=4mm and 4mm]
\node[mkblk,minimum width=13mm] (grid) {Supply\\ (DC bus)};
\node[mkblk,right=of grid,minimum width=13mm] (drv) {Drive\\ (power\\ electronics)};
\node[mkblk,right=of drv,fill=mkCream,draw=mkOrange,minimum width=14mm] (mot) {\textbf{Motor}\\ el.$\rightarrow$mech.};
\node[mkblk,right=of mot,minimum width=13mm] (tr) {Gearbox};
\node[mkblk,right=of tr,minimum width=11mm] (ld) {Robot\\ link};
\draw[mksig,line width=1.4pt,mkOrange] (grid) -- (drv);
\draw[mksig,line width=1.4pt,mkOrange] (drv) -- node[above,mklab]{$v,i$} (mot);
\draw[mksig,line width=1.4pt,mkOrange] (mot) -- node[above,mklab]{$T,\omega$} (tr);
\draw[mksig,line width=1.4pt,mkOrange] (tr) -- (ld);
\node[mkblk,below=7mm of mot,minimum width=13mm,fill=white] (sens) {Encoder};
\node[mkblk,below=7mm of drv,minimum width=13mm,fill=white] (ctl) {Controller};
\draw[mksig] (mot.south) -- (sens.north);
\draw[mksig] (sens) -- (ctl);
\draw[mksig] (ctl) -- (drv);
\node[left=4mm of ctl,mklab,align=center] (ref) {desired\\ angle};
\draw[mksig] (ref) -- (ctl);
\foreach \n/\t in {drv/{heat},mot/{heat},tr/{friction}}{
  \draw[-{Stealth},mkRed,line width=.6pt] (\n.north) -- ++(0,3.5mm) node[above,font=\scriptsize,text=mkRed]{\t};}
\end{tikzpicture}

\vspace{2mm}
\begin{itemize}\setlength\itemsep{1pt}\small
\item Each joint of a 6-axis arm repeats this chain: \hl{6 motors, 6 drives, 6 encoders}.
\item The \term{actuator} is the part that turns energy into motion.
\item Energy flows right; information (angle) flows back to the controller.
\item Every block loses some energy as \hl{heat}.
\end{itemize}
\end{column}
\begin{column}{.29\textwidth}
\centering
\figph{width=\linewidth,height=4.2cm,keepaspectratio}{d4_kuka_robot.jpg}\\
\credit{Oleg Yunakov, CC BY-SA 4.0, via Wikimedia Commons}\\[1mm]
{\scriptsize Industrial arm: one motor + gearbox per joint.}
\end{column}
\end{columns}
\end{frame}

% ---------------------------------------------------------------------
\begin{frame}{Refresher: torque, speed and power}
\begin{columns}[T]
\begin{column}{.42\textwidth}
\centering
\begin{tikzpicture}[scale=.95]
\fill[mkGray] (0,0) circle (0.18);
\draw[mkBlue,line width=3pt] (0,0) -- (3,0);
\draw[-{Stealth},mkOrange,line width=1.4pt] (3,0) -- (3,-1.2) node[below,font=\small]{$F$};
\draw[|-|,mkNavy] (0,0.35) -- node[above,font=\small]{$r$} (3,0.35);
\draw[-{Stealth},mkNavy,line width=1pt] (0.7,-0.5) arc (-35:-145:0.6);
\node[font=\small,text=mkNavy] at (0,-0.95) {$T=F\,r$};
\node[font=\scriptsize,text=mkGray] at (1.5,-1.9) {spanner on a bolt $=$ robot link on a joint};
\end{tikzpicture}

\vspace{2mm}
\small
\renewcommand{\arraystretch}{1.1}
\begin{tabular}{@{}lll@{}}
\toprule
\textbf{Linear} & \textbf{Rotation} & \textbf{Unit}\\
\midrule
force $F$ & torque $T$ & \si{\newton\metre}\\
velocity $v$ & angular speed $\omega$ & \si{\radian\per\second}\\
$P=F\,v$ & $P=T\,\omega$ & \si{\watt}\\
\bottomrule
\end{tabular}
\end{column}
\begin{column}{.55\textwidth}
\small
\textbf{Analogy:} bicycle gears -- the same pedalling power can be ``lots of force, slowly'' or ``little force, fast''.
\begin{eqbox}
$P=T\,\omega$, \qquad $\omega=\dfrac{2\pi n}{60}$ \quad ($n$ in rpm)
\end{eqbox}
\begin{defbox}[Worked example: one robot joint]
\small
\begin{enumerate}\setlength\itemsep{1pt}
\item Joint needs $T=\SI{50}{\newton\metre}$ at $\omega=\SI{1.05}{\radian\per\second}$ (60$^\circ$/s).
\item Power: $P=50\cdot1.05\approx\SI{52}{\watt}$.
\item A motor at \SI{1000}{rpm} ($\SI{105}{\radian\per\second}$) gives this with only $T=52/105\approx\SI{0.5}{\newton\metre}$ -- through a 100:1 gearbox.
\end{enumerate}
\end{defbox}
\hl{Take-away:} power is the product; gears only trade torque for speed.
\end{column}
\end{columns}
\end{frame}

% ---------------------------------------------------------------------
\begin{frame}{Refresher: magnetic field and the force on a wire}
\begin{columns}[T]
\begin{column}{.43\textwidth}
\centering
\begin{tikzpicture}[scale=.95]
\fill[mkRed!70] (0,2.4) rectangle (4,2.9); \node[white,font=\scriptsize\bfseries] at (2,2.65) {N};
\fill[mkMid] (0,-.4) rectangle (4,-.9); \node[white,font=\scriptsize\bfseries] at (2,-.65) {S};
\foreach \x in {0.4,1.0,...,3.8} \draw[-{Stealth[length=1.6mm]},mkSky] (\x,2.3) -- (\x,-.3);
\node[mkSky,font=\scriptsize] at (4.25,1.0) {$\bm{B}$};
\draw[mkGray,line width=1pt] (0.2,0.3) -- (3.9,0.3); \draw[mkGray,line width=1pt] (0.2,1.7) -- (3.9,1.7);
\fill[mkOrange] (2.1,0.2) rectangle (2.35,1.8);
\draw[-{Stealth},mkOrange,line width=1pt] (2.22,0.45) -- (2.22,1.55);
\node[font=\scriptsize,text=mkOrange,right] at (2.35,1.25) {$i$};
\draw[-{Stealth},mkNavy,line width=1.4pt] (2.5,1.0) -- (3.5,1.0) node[right,font=\scriptsize]{$F$};
\draw[|-|,mkNavy] (1.8,0.3) -- node[left,font=\scriptsize]{$l$} (1.8,1.7);
\node[font=\scriptsize] at (2,-1.2) {\color{mkGray}wire on two rails between magnet poles};
\end{tikzpicture}
\end{column}
\begin{column}{.55\textwidth}
\small
\begin{itemize}\setlength\itemsep{2pt}
\item A magnet fills space with a \term{magnetic field} $\bm{B}$ (unit: tesla, \si{\tesla}). Picture it as arrows from N to S.
\item A wire carrying current through that field is \hl{pushed sideways} -- perpendicular to both $i$ and $\bm{B}$.
\item Direction: right-hand rule (fingers along $i$, bend toward $\bm{B}$, thumb $=$ force).
\end{itemize}
\begin{eqbox}
$F = B\,l\,i$ \qquad $F$ [\si{\newton}], $B$ [\si{\tesla}], $l$ [\si{\metre}], $i$ [\si{\ampere}]
\end{eqbox}
\textbf{Example:} $B=\SI{0.5}{\tesla}$, $l=\SI{0.1}{\metre}$, $i=\SI{4}{\ampere}$ $\Rightarrow$ $F=\SI{0.2}{\newton}$.\\[1mm]
\hl{Take-away:} current $+$ magnetic field $=$ force. Every electric motor uses this.
\end{column}
\end{columns}
\end{frame}

% ---------------------------------------------------------------------
\begin{frame}{Motor action and generator action: two halves of every motor}
\begin{columns}[T]
\begin{column}{.47\textwidth}
\small
\textbf{1. Motor action (Lorentz):} current in a field gives force, $F=Bli$.\\[1mm]
\textbf{2. Generator action (Faraday):} a wire \hl{moving} at speed $v$ through the field gets a voltage $e=Blv$, the \term{back-EMF}.\\[1mm]
Both happen at the same time in every motor.\\[2mm]
Bend the wire into a coil ($N$ turns, radius $r$) and both laws use the same constant $k=2NBlr$:
\begin{eqbox}
$T=k_t\,i$, \qquad $e=k_e\,\omega$, \qquad $k_t=k_e$ (SI units)
\end{eqbox}
\end{column}
\begin{column}{.5\textwidth}
\begin{defbox}[Worked example]
\small
Coil: $N=50$, $B=\SI{0.5}{\tesla}$, $l=\SI{0.04}{\metre}$, $r=\SI{0.02}{\metre}$
\begin{enumerate}\setlength\itemsep{1pt}
\item $k=2\cdot50\cdot0.5\cdot0.04\cdot0.02=\SI{0.04}{\newton\metre\per\ampere}$
\item At $i=\SI{5}{\ampere}$: $T=0.04\cdot5=\SI{0.2}{\newton\metre}$
\item At $\omega=\SI{300}{\radian\per\second}$: $e=0.04\cdot300=\SI{12}{\volt}$
\end{enumerate}
\end{defbox}
\begin{keybox}[Take-away]
\small
\term{Torque constant} $k_t$ (\si{\newton\metre\per\ampere}) and \term{back-EMF constant} $k_e$ (\si{\volt\second\per\radian}) are the same number: \hl{current makes torque, speed makes voltage}.
\end{keybox}
\end{column}
\end{columns}
\end{frame}

% ---------------------------------------------------------------------
\begin{frame}{Reluctance force: iron is pulled into the field}
\begin{columns}[T]
\begin{column}{.40\textwidth}
\centering
\begin{tikzpicture}[scale=.8]
\fill[mkGray!50] (0,0) rectangle (0.6,3); \fill[mkGray!50] (0,0) rectangle (2.4,0.6); \fill[mkGray!50] (0,2.4) rectangle (2.4,3);
\foreach \y in {0.9,1.2,...,2.1} \draw[mkOrange,line width=1.6pt] (-0.15,\y) -- (0.75,\y);
\node[font=\scriptsize,text=mkOrange] at (-0.5,1.5) {$N,i$};
\begin{scope}[shift={(2.85,1.5)},rotate=30]
\fill[mkMid] (-0.25,-1.1) rectangle (0.25,1.1);
\end{scope}
\draw[-{Stealth},mkNavy,line width=.8pt] (3.9,0.5) arc (-40:40:1.3) node[right,font=\scriptsize]{$T$};
\node[font=\scriptsize] at (2.85,-0.3) {iron rotor at angle $\theta$};
\end{tikzpicture}
\begin{tikzpicture}
\begin{axis}[mkplot,width=5.2cm,height=3.4cm,xmin=0,xmax=180,xtick={0,45,90,135,180},
  xlabel={$\theta$ (deg)},ymin=-1.1,ymax=2.1,legend pos=south west,legend columns=2]
\addplot[mkBlue,thick,domain=0:180,samples=60]{1+0.8*cos(2*x)}; \addlegendentry{$L(\theta)/L_0$}
\addplot[mkOrange,thick,domain=0:180,samples=60]{-0.8*sin(2*x)}; \addlegendentry{$T/T_0$}
\end{axis}
\end{tikzpicture}
\end{column}
\begin{column}{.57\textwidth}
\small
\textbf{Analogy:} a fridge magnet snaps onto the steel door -- the iron moves to where the magnetic path is \hl{easiest}.
\begin{itemize}\setlength\itemsep{2pt}
\item No magnet needed on the rotor: a plain iron bar turns to line up with the coil.
\item Torque pushes towards the position of \hl{highest inductance} $L$.
\item Torque grows with $i^2$: current direction does not matter.
\item Used in solenoids, \term{stepper motors} (Session 1.4) and switched-reluctance motors.
\end{itemize}
\begin{eqbox}
$T=\tfrac12\,i^2\,\dfrac{dL}{d\theta}$ \quad ($L$ in \si{\henry}, $\theta$ in rad)
\end{eqbox}
\gofurther{the co-energy method $T=\partial W'/\partial\theta$ also covers saturated iron [Fitzgerald et al., 2003, ch.\,3].}
\end{column}
\end{columns}
\end{frame}

% ---------------------------------------------------------------------
\begin{frame}{The actuator landscape: three ways to move a robot}
\centering
\begin{tikzpicture}[lbl/.style={font=\scriptsize,text=mkNavy,align=center}]
% electric
\begin{scope}[shift={(0,0)}]
\node[font=\small\bfseries,text=mkBlue] at (1.6,2.3) {Electric};
\draw[mkBlue,fill=mkPale,rounded corners=2pt] (0,0) rectangle (2.2,1.5);
\foreach \x in {0.3,0.6,...,1.9} \draw[mkOrange,line width=1pt] (\x,0.15) -- (\x,1.35);
\fill[mkGray] (2.2,0.65) rectangle (3.3,0.85);
\draw[mkNavy,line width=.8pt] (-0.7,0.4) -- (0,0.4) (-0.7,1.1) -- (0,1.1);
\node[lbl] at (-0.7,1.45) {cable};
\draw[-{Stealth},mkNavy,line width=1pt] (3.0,1.25) arc (90:-200:0.33);
\node[lbl] at (1.4,-0.5) {motor: $v,i\;\rightarrow\;T,\omega$};
\end{scope}
% pneumatic
\begin{scope}[shift={(5.2,0)}]
\node[font=\small\bfseries,text=mkBlue] at (1.4,2.3) {Pneumatic};
\draw[mkBlue,fill=mkLight] (0,0.2) rectangle (2.6,1.3);
\fill[mkSky!40] (0.02,0.22) rectangle (1.2,1.28);
\fill[mkGray] (1.2,0.2) rectangle (1.4,1.3);
\fill[mkGray] (1.4,0.65) rectangle (3.4,0.85);
\draw[mkSky,line width=1.5pt] (0.3,1.3) -- (0.3,1.8) -- (-0.6,1.8);
\node[lbl] at (-0.4,2.05) {air 6 bar};
\draw[-{Stealth},mkNavy,line width=1pt] (2.8,1.1) -- (3.5,1.1);
\node[lbl] at (1.4,-0.5) {cylinder: $p\;\rightarrow\;F$};
\end{scope}
% hydraulic
\begin{scope}[shift={(10.2,0)}]
\node[font=\small\bfseries,text=mkBlue] at (1.4,2.3) {Hydraulic};
\draw[mkBlue,fill=mkLight,line width=1.2pt] (0,0.2) rectangle (2.6,1.3);
\fill[mkAmber!60] (0.02,0.22) rectangle (1.2,1.28);
\fill[mkGray] (1.2,0.2) rectangle (1.4,1.3);
\fill[mkGray] (1.4,0.62) rectangle (3.4,0.88);
\draw[mkAmber,line width=1.5pt] (0.3,1.3) -- (0.3,1.8) -- (-0.6,1.8);
\node[lbl] at (-0.4,2.05) {oil 200 bar};
\draw[-{Stealth},mkNavy,line width=1.6pt] (2.8,1.1) -- (3.5,1.1);
\node[lbl] at (1.4,-0.5) {cylinder: $p\;\rightarrow\;F$ (large)};
\end{scope}
\end{tikzpicture}

\vspace{3mm}
\begin{columns}[T]
\begin{column}{.5\textwidth}
\small
\begin{itemize}\setlength\itemsep{2pt}
\item \term{Electric}: energy arrives by cable; easiest to \hl{control precisely}.
\item \term{Pneumatic}: compressed air pushes a piston; cheap, fast, \hl{springy}.
\item \term{Hydraulic}: oil at high pressure; \hl{huge force} from a small cylinder.
\end{itemize}
\end{column}
\begin{column}{.46\textwidth}
\begin{notebox}
\small
Today's industrial and collaborative robots are almost all \hl{electric}. Pneumatics live in grippers and simple pick-and-place; hydraulics in heavy and legged machines.
\end{notebox}
\end{column}
\end{columns}
\end{frame}

% ---------------------------------------------------------------------
\begin{frame}{Fluid power in numbers: force from pressure}
\begin{columns}[T]
\begin{column}{.5\textwidth}
\centering
\begin{tikzpicture}[scale=.75]
\draw[mkBlue,line width=1pt,fill=mkLight] (0,0) rectangle (4,1.6);
\fill[mkSky!40] (0.03,0.03) rectangle (1.8,1.57);
\fill[mkGray] (1.8,0) rectangle (2.1,1.6);
\fill[mkGray] (2.1,0.65) rectangle (5.2,0.95);
\foreach \y in {0.25,0.55,...,1.45} \draw[-{Stealth},mkOrange,line width=.7pt] (1.2,\y) -- (1.75,\y);
\node[font=\small,text=mkOrange] at (0.7,0.8) {$p$};
\draw[|-|,mkNavy] (2.4,0) -- node[right,font=\small]{$d$} (2.4,1.6);
\draw[-{Stealth},mkNavy,line width=1.4pt] (5.25,0.8) -- (6.2,0.8) node[right,font=\small]{$F$};
\node[font=\scriptsize,text=mkGray] at (2.5,-0.4) {piston area $A=\pi d^2/4$};
\end{tikzpicture}

\vspace{2mm}
\figph{width=.9\linewidth,height=2.1cm,keepaspectratio}{r11_pneumatic.jpg}\\
\credit{Jonashtand, CC BY-SA 4.0, via Wikimedia Commons}
\end{column}
\begin{column}{.47\textwidth}
\small
\begin{eqbox}
$F=p\,A$ \quad (\si{\newton}, \si{\pascal}, \si{\metre\squared}); \; \SI{1}{bar} $=\SI{e5}{\pascal}$
\end{eqbox}
\begin{defbox}[Worked example: bore 40 mm]
\small
\begin{enumerate}\setlength\itemsep{1pt}
\item $A=\pi\cdot0.04^2/4=\SI{1.26e-3}{\metre\squared}$
\item Pneumatic, \SI{6}{bar}: $F=\num{6e5}\cdot\num{1.26e-3}\approx\SI{750}{\newton}$
\item Hydraulic, \SI{200}{bar}: $F\approx\SI{25}{\kilo\newton}$ -- \hl{33$\times$ more}
\end{enumerate}
\end{defbox}
Air is springy: \hl{hard to stop mid-stroke}. Oil is stiff, but needs pumps and hoses.
\end{column}
\end{columns}
\end{frame}

% ---------------------------------------------------------------------
\begin{frame}{Comparing actuator technologies for robots}
\begin{columns}[T]
\begin{column}{.68\textwidth}
\small
\renewcommand{\arraystretch}{1.15}
\begin{tabularx}{\linewidth}{@{}l*{3}{>{\centering\arraybackslash}X}@{}}
\toprule
 & \textbf{Electric} & \textbf{Pneumatic} & \textbf{Hydraulic}\\
\midrule
Force per size & medium--high & low & \textcolor{mkGreen}{\textbf{very high}}\\
Positioning accuracy & \textcolor{mkGreen}{\textbf{very good}} & \textcolor{mkRed}{poor (2 end stops)} & good\\
Control & \textcolor{mkGreen}{\textbf{easy, digital}} & on/off valves & servo valves\\
Efficiency & \textcolor{mkGreen}{\textbf{80--95\,\%}} & \textcolor{mkRed}{10--20\,\%} & 40--60\,\%\\
Cost / infrastructure & medium & \textcolor{mkGreen}{\textbf{low}} & \textcolor{mkRed}{pump, hoses, leaks}\\
Typical robot use & arm joints, AMRs & grippers, suction & heavy, legged robots\\
\bottomrule
\end{tabularx}
\end{column}
\begin{column}{.29\textwidth}
\centering
\figph{width=\linewidth,height=2.6cm,keepaspectratio}{r11_hydraulic.jpg}\\
\credit{S.J. de Waard, CC BY 2.5, via Wikimedia Commons}\\[1mm]
{\scriptsize Hydraulic cylinders: large forces in heavy machines.}
\end{column}
\end{columns}

\vspace{1mm}
\begin{keybox}[Why robots went electric]
\small
Clean, efficient and easy to control with a computer and an encoder; gearboxes (Day 2) fix their weakness of low torque.
\end{keybox}
\end{frame}

% ---------------------------------------------------------------------
\begin{frame}{Family tree of electromagnetic actuators}
\centering
\resizebox{\textwidth}{!}{\begin{tikzpicture}[
  lvl1/.style={mkblk,fill=mkBlue,text=white,minimum width=24mm,font=\scriptsize\bfseries},
  lvl2/.style={mkblk,minimum width=20mm,text width=20mm},
  lvl3/.style={font=\scriptsize,text=mkNavy,align=left,text width=22mm},
  edge/.style={mkGray,line width=.5pt}]
\node[mkblk,fill=mkNavy,text=white,font=\small\bfseries,minimum width=40mm] (root) {Robot actuators};
\node[lvl1,below left=6mm and 26mm of root] (em) {Electromagnetic};
\node[lvl1,below right=6mm and 26mm of root] (oth) {Other principles};
\draw[edge] (root) -- (em); \draw[edge] (root) -- (oth);
\node[lvl2,below=7mm of em,xshift=-36mm] (dc) {DC (brushed)};
\node[lvl2,right=2mm of dc,fill=mkCream,draw=mkOrange] (pm) {PM brushless\\ BLDC / PMSM};
\node[lvl2,right=2mm of pm] (im) {Induction};
\node[lvl2,right=2mm of im] (rel) {Stepper / reluctance};
\node[lvl2,right=2mm of rel] (lin) {Linear motor, solenoid};
\foreach \n in {dc,pm,im,rel,lin} \draw[edge] (em.south) -- (\n.north);
\node[lvl3,below=1mm of dc] {toy and hobby robots, small grippers};
\node[lvl3,below=1mm of pm] {\textbf{arm joints}, cobots, AMR wheels, drones};
\node[lvl3,below=1mm of im] {conveyors, pumps, big spindles};
\node[lvl3,below=1mm of rel] {3D printers, small gantries};
\node[lvl3,below=1mm of lin] {fast linear axes, valves};
\node[lvl2,right=8mm of lin] (fl) {Fluid power};
\node[lvl2,right=2mm of fl] (sm) {Smart materials};
\draw[edge] (oth.south) -- (fl.north); \draw[edge] (oth.south) -- (sm.north);
\node[lvl3,below=1mm of fl] {pneumatic grippers, hydraulic legs};
\node[lvl3,below=1mm of sm] {piezo, shape-memory alloy, soft actuators};
\end{tikzpicture}}

\vspace{4mm}
\begin{columns}[T]
\begin{column}{.55\textwidth}
\small
\begin{itemize}\setlength\itemsep{2pt}
\item Today: the \hl{rotary electromagnetic} branch (Sessions 1.2--1.4).
\item Day 2: gearboxes, linear actuators, grippers, sizing.
\end{itemize}
\end{column}
\begin{column}{.41\textwidth}
\gofurther{soft robots use pneumatic ``muscles'' and electro-active polymers -- high compliance, low precision.}
\end{column}
\end{columns}
\end{frame}

% ---------------------------------------------------------------------
\begin{frame}{Session 1.1 -- Check your understanding}
\begin{columns}[T]
\begin{column}{.62\textwidth}
\begin{keybox}[Check your understanding]
\begin{enumerate}\setlength\itemsep{6pt}
\item A motor has $k_t=\SI{0.1}{\newton\metre\per\ampere}$. You double the current from \SI{2}{\ampere} to \SI{4}{\ampere}. What happens to the torque?
\item A joint delivers \SI{2}{\newton\metre} at \SI{1500}{rpm}. Roughly how much mechanical power is that?
\item Why is a pneumatic cylinder a poor choice when a robot must stop accurately halfway along its stroke?
\end{enumerate}
\end{keybox}
\end{column}
\begin{column}{.34\textwidth}
\small
\textbf{How to work:}
\begin{itemize}
\item 2 minutes alone, then compare with your neighbour.
\item Write the formula first, then the numbers with units.
\item Hint: $T=k_t i$, $P=T\omega$, $F=pA$.
\end{itemize}
\end{column}
\end{columns}
\end{frame}

% ---------------------------------------------------------------------
\begin{frame}{Session 1.1 -- Summary and key terms}
\begin{columns}[T]
\begin{column}{.55\textwidth}
\small
\begin{itemize}
\item A motor turns $v\,i$ into $T\,\omega$ through a magnetic field -- and back again when braking.
\item Current in a field gives force; motion in a field gives voltage.
\item Iron alone is also pulled into a field (reluctance).
\item Electric actuators dominate robots; pneumatics for grippers, hydraulics for heavy loads.
\end{itemize}
\begin{keybox}[Remember]
$P=T\omega$, \quad $T=k_t\,i$, \quad $e=k_e\,\omega$, \quad $F=p\,A$
\end{keybox}
\end{column}
\begin{column}{.42\textwidth}
\begin{defbox}[Key terms]
\small
\term{Actuator}: device that turns energy into motion.\\
\term{Torque $T$}: turning effect of a force, \si{\newton\metre}.\\
\term{Back-EMF}: voltage a spinning motor generates.\\
\term{Torque constant $k_t$}: torque per ampere.\\
\term{Reluctance}: ``difficulty'' of the magnetic path.\\
\term{Efficiency $\eta$}: power out / power in.
\end{defbox}
\end{column}
\end{columns}
\end{frame}

% =====================================================================
\section{1.2 DC, BLDC and PMSM motors}
% =====================================================================

\begin{frame}{Anatomy of a robot servo motor}
\begin{columns}[T]
\begin{column}{.47\textwidth}
\centering
\begin{tikzpicture}[scale=.72,lbl/.style={font=\scriptsize,text=mkNavy},ld/.style={mkGray,line width=.3pt}]
\draw[mkGray,line width=.8pt,fill=mkLight] (0,0) rectangle (6.4,3.4);
\foreach \x in {0.6,1.0,...,4.6} {\fill[mkGray!40] (\x,3.4) rectangle (\x+0.15,3.6); \fill[mkGray!40] (\x,-0.2) rectangle (\x+0.15,0);}
\fill[mkGray!60] (-0.35,-0.35) rectangle (0,3.75);
\fill[mkSky!60] (1.0,2.3) rectangle (4.4,3.2); \fill[mkSky!60] (1.0,0.2) rectangle (4.4,1.1);
\foreach \y in {2.3,0.2} {\fill[mkOrange!80] (0.6,\y+0.1) rectangle (1.0,\y+0.8); \fill[mkOrange!80] (4.4,\y+0.1) rectangle (4.8,\y+0.8);}
\fill[mkGray!50] (1.0,1.25) rectangle (4.4,2.15);
\fill[mkRed!70] (1.0,2.0) rectangle (4.4,2.15); \fill[mkMid] (1.0,1.25) rectangle (4.4,1.4);
\fill[mkGray] (-1.3,1.55) rectangle (5.9,1.85);
\foreach \x in {0.25,5.0} {\draw[mkNavy,fill=white] (\x,1.85) rectangle (\x+0.35,2.2); \draw[mkNavy] (\x,1.85) -- (\x+0.35,2.2) (\x,2.2) -- (\x+0.35,1.85);
  \draw[mkNavy,fill=white] (\x,1.2) rectangle (\x+0.35,1.55); \draw[mkNavy] (\x,1.2) -- (\x+0.35,1.55) (\x,1.55) -- (\x+0.35,1.2);}
\fill[mkAmber!70] (5.45,1.0) rectangle (5.7,2.4);
\fill[mkNavy] (5.9,1.1) rectangle (5.98,2.3); \fill[mkGreen!70] (6.05,2.0) rectangle (6.3,2.3);
\fill[mkGray!70] (2.6,3.6) rectangle (3.4,4.1);
\draw[ld] (2.7,2.75) -- (2.7,4.5) node[above,lbl]{stator + winding};
\draw[ld] (3.4,3.85) -- (5.0,4.5) node[above,lbl]{connector};
\draw[ld] (2.0,2.07) -- (1.6,-0.9) node[below,lbl]{magnet rotor};
\draw[ld] (0.42,1.3) -- (0.1,-0.9) node[below,lbl]{bearing};
\draw[ld] (-0.9,1.7) -- (-1.3,-0.9) node[below,lbl]{shaft};
\draw[ld] (5.57,1.1) -- (4.4,-0.9) node[below,lbl]{brake};
\draw[ld] (6.1,2.1) -- (6.4,-1.6) node[below,lbl]{encoder};
\draw[ld] (3.0,0.9) -- (3.0,-1.6) node[below,lbl]{air gap};
\end{tikzpicture}\\
{\scriptsize\color{mkGray} Schematic cut through a PM servo motor (not to scale).}
\end{column}
\begin{column}{.5\textwidth}
\small
\begin{itemize}\setlength\itemsep{2pt}
\item \term{Stator}: fixed outer part -- iron stack with copper windings. Gets hot, so it touches the housing.
\item \term{Rotor}: turning part -- shaft with magnets. Long and slim, so it is light to accelerate.
\item \term{Air gap}: thin gap ($<\SI{1}{\milli\metre}$) where torque is made.
\item \term{Encoder}: measures the rotor angle (Week 2).
\item \term{Brake}: holds the arm when power is off -- vital on a robot, or the arm falls.
\end{itemize}
\end{column}
\end{columns}
\end{frame}

% ---------------------------------------------------------------------
\begin{frame}{Refresher: a magnet follows a rotating field}
\begin{columns}[T]
\begin{column}{.5\textwidth}
\centering
\begin{tikzpicture}[scale=.62]
\foreach \cx/\ang/\lab in {0/90/{A on},3.3/30/{A + B on},6.6/-30/{B on}}{
\begin{scope}[shift={(\cx,0)}]
\draw[mkGray,line width=3pt] (0,0) circle (1.25);
\foreach \a/\n in {90/A,-30/B,210/C}{
  \fill[mkOrange!80,rotate=\a] (1.0,-0.25) rectangle (1.2,0.25);
  \node[font=\tiny,text=mkNavy] at (\a:1.55) {\n};}
\draw[-{Stealth},mkSky,line width=2pt] (0,0) -- (\ang:0.95);
\begin{scope}[rotate=\ang]
\fill[mkRed!70] (0,-0.13) rectangle (0.6,0.13);
\fill[mkMid] (-0.6,-0.13) rectangle (0,0.13);
\end{scope}
\fill[mkGray] (0,0) circle (0.06);
\node[font=\scriptsize,text=mkNavy] at (0,-1.9) {\lab};
\end{scope}}
\end{tikzpicture}

\vspace{2mm}
{\scriptsize\color{mkGray} Stator field (blue arrow) moves step by step; the magnet rotor (red = N) follows.}
\end{column}
\begin{column}{.47\textwidth}
\small
\textbf{Analogy:} a compass needle next to a coil -- switch coils around the circle and the needle chases the field like a donkey chasing a carrot.
\begin{itemize}\setlength\itemsep{2pt}
\item Three coils A, B, C, $120^\circ$ apart, make a field we can \hl{point anywhere} by choosing their currents.
\item The rotor magnet always tries to line up with it.
\item Torque is largest when field and magnet are \hl{$90^\circ$ apart}, zero when aligned.
\end{itemize}
\begin{eqbox}
$T\propto B_\text{stator}\,B_\text{rotor}\,\sin\delta$ \quad ($\delta$: angle between them)
\end{eqbox}
\hl{Take-away:} commutation = keep the field $90^\circ$ ahead of the rotor.
\end{column}
\end{columns}
\end{frame}

% ---------------------------------------------------------------------
\begin{frame}{Brushed DC motor: construction}
\begin{columns}[T]
\begin{column}{.36\textwidth}
\centering
\begin{tikzpicture}[scale=.85]
\fill[mkGray!40] (0,0) circle (2.0); \fill[white] (0,0) circle (1.65);
\fill[mkRed!70] (100:1.65) arc (100:260:1.65) -- (260:1.35) arc (260:100:1.35) -- cycle;
\fill[mkMid] (-80:1.65) arc (-80:80:1.65) -- (80:1.35) arc (80:-80:1.35) -- cycle;
\node[white,font=\scriptsize\bfseries] at (180:1.5) {N}; \node[white,font=\scriptsize\bfseries] at (0:1.5) {S};
\fill[mkSky!50] (0,0) circle (1.22);
\foreach \a in {0,30,...,330}{ \fill[white] (\a:1.05) circle (0.1);}
\foreach \a in {120,150,180,210,240}{ \draw[mkOrange,line width=.7pt] (\a:1.05) circle (0.09); \fill[mkOrange] (\a:1.05) circle (0.025);}
\foreach \a in {-60,-30,0,30,60}{ \draw[mkOrange,line width=.7pt] (\a:1.05) circle (0.09);
   \draw[mkOrange,line width=.6pt] ($(\a:1.05)+(-0.06,-0.06)$) -- ($(\a:1.05)+(0.06,0.06)$) ($(\a:1.05)+(-0.06,0.06)$) -- ($(\a:1.05)+(0.06,-0.06)$);}
\fill[mkAmber] (0,0) circle (0.42); \foreach \a in {0,45,...,315} \draw[white] (0,0) -- (\a:0.42);
\fill[mkNavy] (-0.42,-0.1) rectangle (-0.7,0.1); \fill[mkNavy] (0.42,-0.1) rectangle (0.7,0.1);
\fill[mkGray] (0,0) circle (0.12);
\draw[-{Stealth},mkNavy,line width=1pt] (60:2.25) arc (60:120:2.25) node[left,font=\scriptsize]{$T$};
\node[font=\scriptsize,text=mkGray] at (0,-2.35) {magnets, armature, commutator};
\end{tikzpicture}
\end{column}
\begin{column}{.27\textwidth}
\centering
\figph{width=\linewidth,height=3.6cm,keepaspectratio}{d1_brushed_dc.jpg}\\
\credit{Jfer91, CC BY-SA 3.0, via Wikimedia Commons}\\[1mm]
{\scriptsize Opened DC motor: magnets in the can, armature and commutator on the shaft.}
\end{column}
\begin{column}{.35\textwidth}
\small
\begin{itemize}\setlength\itemsep{2pt}
\item \term{Stator magnets} make the field.
\item \term{Armature}: the coils on the \hl{rotor} that carry the current.
\item \term{Commutator + brushes}: a rotating switch that feeds the coils.
\item Simple: connect a battery and it spins; speed $\approx$ proportional to voltage.
\item Found in toy robots, small grippers, car wipers.
\end{itemize}
\end{column}
\end{columns}
\end{frame}

% ---------------------------------------------------------------------
\begin{frame}{How the commutator keeps the torque in one direction}
\begin{columns}[T]
\begin{column}{.52\textwidth}
\centering
\begin{tikzpicture}[scale=.62]
\foreach \cx/\ang/\lab/\flip in {0/0/{$0^\circ$: max}/0,3.6/80/{$90^\circ$: switch}/0,7.2/180/{$180^\circ$: max again}/1}{
\begin{scope}[shift={(\cx,0)}]
\fill[mkRed!70] (-1.6,-1.0) rectangle (-1.35,1.0); \node[white,font=\tiny\bfseries] at (-1.475,0) {N};
\fill[mkMid] (1.35,-1.0) rectangle (1.6,1.0); \node[white,font=\tiny\bfseries] at (1.475,0) {S};
\draw[mkOrange,line width=1.4pt,rotate=\ang] (0,-0.9) -- (0,0.9);
\fill[mkOrange,rotate=\ang] (0,0.9) circle (0.12);
\fill[mkOrange,rotate=\ang] (0,-0.9) circle (0.12);
\ifnum\flip=0
 \fill[white,rotate=\ang] (0,0.9) circle (0.05); \draw[white,line width=.6pt,rotate=\ang] (-0.07,-0.97) -- (0.07,-0.83) (-0.07,-0.83) -- (0.07,-0.97);
\else
 \fill[white,rotate=\ang] (0,-0.9) circle (0.05); \draw[white,line width=.6pt,rotate=\ang] (-0.07,0.97) -- (0.07,0.83) (-0.07,0.83) -- (0.07,0.97);
\fi
\node[font=\scriptsize,text=mkNavy,align=center] at (0,-1.5) {\lab};
\end{scope}}
\end{tikzpicture}

\begin{tikzpicture}
\begin{axis}[mkplot,width=7.0cm,height=3.3cm,xmin=0,xmax=360,ymin=0,ymax=1.15,
  xtick={0,90,180,270,360},xlabel={rotor angle (deg)},ylabel={$T/T_\text{max}$},legend pos=south east,legend columns=2]
\addplot[mkSky,thick,domain=0:360,samples=181]{abs(cos(x))}; \addlegendentry{1 coil}
\addplot[mkOrange,thick,domain=0:360,samples=361]{max(abs(cos(x)),abs(cos(x+60)),abs(cos(x+120)))}; \addlegendentry{3 coils}
\end{axis}
\end{tikzpicture}
\end{column}
\begin{column}{.45\textwidth}
\small
\begin{itemize}\setlength\itemsep{2pt}
\item Without switching, the coil would rock back and forth: torque changes sign every half turn.
\item The \term{commutator} reverses the coil current exactly when the coil passes the ``dead'' position.
\item Result: torque always pushes \hl{the same way} (rectified $|\cos|$).
\item More coils and segments $\Rightarrow$ smaller dips $=$ less \term{torque ripple}.
\end{itemize}
\begin{keybox}[Weak points of brushes]
\small
Brush wear and sparks, heat trapped in the rotor, speed limit. \hl{Idea:} swap roles -- magnets on the rotor, coils on the stator, switches made of transistors.
\end{keybox}
\end{column}
\end{columns}
\end{frame}

% ---------------------------------------------------------------------
\begin{frame}{Brushless DC (BLDC): the ``inside-out'' DC motor}
\begin{columns}[T]
\begin{column}{.45\textwidth}
\centering
\begin{tikzpicture}[scale=.8]
\fill[mkGray!40] (0,0) circle (2.0); \fill[white] (0,0) circle (1.48);
\foreach \a/\n in {90/A,30/C',-30/B,-90/A',-150/C,150/B'}{
  \fill[mkGray!40,rotate=\a] (1.2,-0.18) rectangle (1.5,0.18);
  \fill[mkOrange!70,rotate=\a] (1.25,-0.3) rectangle (1.45,0.3);
  \node[font=\tiny\bfseries,text=mkNavy] at (\a:1.75) {\n};}
\fill[mkSky!50] (0,0) circle (0.95);
\fill[mkRed!70] (0,0) -- (-0.95,0) arc (180:0:0.95) -- cycle;
\fill[mkMid] (0,0) -- (0.95,0) arc (0:-180:0.95) -- cycle;
\node[white,font=\scriptsize\bfseries] at (0,0.5) {N}; \node[white,font=\scriptsize\bfseries] at (0,-0.5) {S};
\fill[mkGray] (0,0) circle (0.13);
\foreach \a/\h in {0/H1,120/H2,240/H3} \node[draw=mkGreen,fill=mkGreen!30,font=\tiny,inner sep=1pt] at (\a+60:1.12) {\h};
\node[font=\scriptsize,text=mkGray] at (0,-2.35) {3 phases A, B, C + Hall sensors H1--H3};
\end{tikzpicture}

\vspace{1mm}
\figph{width=.62\linewidth,height=2.6cm,keepaspectratio}{d4_bldc_stator.jpg}\\
\credit{Medvedev, CC BY-SA 3.0, via Wikimedia Commons}
\end{column}
\begin{column}{.52\textwidth}
\small
\begin{itemize}\setlength\itemsep{2pt}
\item \hl{Magnets on the rotor}, three-phase coils on the stator (photo) $\Rightarrow$ no brushes, no wear.
\item Heat is made in the stator, right next to the housing $\Rightarrow$ \hl{easy cooling}.
\item The drive must switch the coils itself, so it needs to know the rotor angle: three cheap \term{Hall sensors}.
\item This is \term{electronic commutation}: transistors replace the commutator.
\end{itemize}
\begin{notebox}
\small
Same physics as the brushed motor ($T=k_ti$, $e=k_e\omega$) -- only the switch moved from copper segments into the drive.
\end{notebox}
\end{column}
\end{columns}
\end{frame}

% ---------------------------------------------------------------------
\begin{frame}{BLDC six-step commutation}
\begin{columns}[T]
\begin{column}{.55\textwidth}
\centering
\begin{tikzpicture}
\begin{groupplot}[group style={group size=1 by 2,vertical sep=4mm,xticklabels at=edge bottom},
  mkplot,width=7.4cm,height=2.8cm,xmin=0,xmax=360,xtick={0,60,...,360},ymin=-1.3,ymax=1.3,ytick={-1,0,1}]
\nextgroupplot[ylabel={$e_a$}]
\addplot[mkSky,thick] coordinates {(0,0) (30,1) (150,1) (210,-1) (330,-1) (360,0)};
\node[font=\tiny,text=mkGray] at (axis cs:90,0.5) {back-EMF of phase A (trapezoid)};
\nextgroupplot[ylabel={$i_a$},xlabel={electrical angle (deg)}]
\addplot[mkOrange,thick,const plot] coordinates {(0,0) (30,1) (150,0) (210,-1) (330,0) (360,0)};
\node[font=\tiny,text=mkGray] at (axis cs:270,0.5) {block current};
\end{groupplot}
\end{tikzpicture}
\end{column}
\begin{column}{.42\textwidth}
\small
\renewcommand{\arraystretch}{1.05}
\begin{tabular}{@{}ccc@{}}
\toprule
\textbf{Hall H1 H2 H3} & \textbf{$+$ phase} & \textbf{$-$ phase}\\
\midrule
1\,0\,1 & A & B\\
1\,0\,0 & A & C\\
1\,1\,0 & B & C\\
0\,1\,0 & B & A\\
0\,1\,1 & C & A\\
0\,0\,1 & C & B\\
\bottomrule
\end{tabular}\\[1mm]
{\scriptsize\color{mkGray}Example sequence; the exact table depends on sensor placement.}
\begin{itemize}\setlength\itemsep{1pt}
\item At any moment \hl{two phases conduct}, one rests.
\item 6 steps per electrical turn, each $60^\circ$.
\item Simple and cheap, but a small \term{torque ripple} at every step.
\end{itemize}
\end{column}
\end{columns}
\end{frame}

% ---------------------------------------------------------------------
\begin{frame}{PMSM: smooth sine currents -- the robot servo motor}
\begin{columns}[T]
\begin{column}{.48\textwidth}
\centering
\begin{tikzpicture}[scale=.62]
\foreach \cx/\name in {0/SPM: magnets on surface, 4.6/IPM: magnets buried}{
\begin{scope}[shift={(\cx,0)}]
 \fill[mkGray!40] (0,0) circle (2.0); \fill[white] (0,0) circle (1.48);
 \foreach \a in {0,30,...,330}{ \fill[mkGray!40,rotate=\a] (1.4,-0.13) rectangle (1.55,0.13);}
 \foreach \a in {15,45,...,345}{ \fill[mkOrange!60,rotate=\a] (1.55,-0.1) rectangle (1.85,0.1);}
 \fill[mkSky!50] (0,0) circle (1.1);
 \fill[mkGray] (0,0) circle (0.25);
 \node[font=\scriptsize,text=mkNavy] at (0,-2.45) {\name};
\end{scope}}
\foreach \a/\c in {0/mkRed!70,90/mkMid,180/mkRed!70,270/mkMid}{
  \fill[\c] (\a-38:1.32) arc (\a-38:\a+38:1.32) -- (\a+38:1.1) arc (\a+38:\a-38:1.1) -- cycle;}
\begin{scope}[shift={(4.6,0)}]
\foreach \a/\c in {0/mkRed!70,90/mkMid,180/mkRed!70,270/mkMid}{
  \begin{scope}[rotate=\a]
  \fill[\c,rotate around={25:(0.75,0.25)}] (0.45,0.2) rectangle (1.0,0.32);
  \fill[\c,rotate around={-25:(0.75,-0.25)}] (0.45,-0.32) rectangle (1.0,-0.2);
  \end{scope}}
\end{scope}
\end{tikzpicture}\\[1mm]
\begin{tikzpicture}
\begin{axis}[mkplot,width=6.4cm,height=3.0cm,xmin=0,xmax=360,ymin=-1.2,ymax=1.2,xtick={0,120,240,360},
xlabel={electrical angle (deg)},ylabel={$i_a,i_b,i_c$},ytick={-1,0,1}]
\addplot[mkBlue,thick,domain=0:360,samples=90]{sin(x)};
\addplot[mkOrange,thick,domain=0:360,samples=90]{sin(x-120)};
\addplot[mkSky,thick,domain=0:360,samples=90]{sin(x+120)};
\end{axis}
\end{tikzpicture}
\end{column}
\begin{column}{.5\textwidth}
\small
\begin{itemize}\setlength\itemsep{2pt}
\item \term{PMSM} (permanent-magnet synchronous motor): same build as BLDC, but sine-shaped back-EMF.
\item An \hl{encoder} gives the exact angle; the drive feeds three smooth \hl{sine currents} $120^\circ$ apart.
\item The field turns smoothly $\Rightarrow$ almost \hl{ripple-free torque}.
\item Industrial ``AC servo motors'' and almost every robot arm joint use a PMSM.
\end{itemize}
\figph{width=.4\linewidth,height=1.6cm,keepaspectratio}{d1_bldc.jpg}
\credit{Sebastian Koppehel, CC BY 3.0, via Wikimedia Commons}\\
\gofurther{field-oriented control computes these sine currents from the rotor angle: $T=\tfrac32\,p\,\psi_\text{PM}\,i_q$ (W1D5).}
\end{column}
\end{columns}
\end{frame}

% ---------------------------------------------------------------------
\begin{frame}{Torque ripple: six-step versus sinusoidal}
\begin{columns}[T]
\begin{column}{.55\textwidth}
\begin{tikzpicture}
\begin{axis}[mkplot,width=7.6cm,height=4.8cm,xmin=0,xmax=360,ymin=0.7,ymax=1.1,xtick={0,60,...,360},
  xlabel={electrical angle (deg)},ylabel={$T/T_\text{avg}$},legend pos=south east]
\addplot[mkSky,thick,domain=0:360,samples=361]{ (1/0.955)*cos( mod(x,60)-30 ) }; \addlegendentry{BLDC, sine back-EMF, block current}
\addplot[mkOrange,thick,domain=0:360,samples=2]{1}; \addlegendentry{PMSM, sine current}
\end{axis}
\end{tikzpicture}\\
{\scriptsize\color{mkGray} Ideal case. With block currents on a sine-EMF motor the dips reach $\approx 13\,\%$ every $60^\circ$.}
\end{column}
\begin{column}{.42\textwidth}
\small
\textbf{Analogy:} a bicycle pedalled with jerky pushes vs.\ smooth circles -- same average power, different feel.
\begin{itemize}\setlength\itemsep{2pt}
\item \term{Torque ripple}: periodic wobble on top of the mean torque.
\item In a robot it shows up as \hl{vibration} at the tool tip and poor slow-speed motion.
\item Sine commutation with an accurate encoder removes most of it.
\end{itemize}
\begin{notebox}
\small
Drones and fans: BLDC six-step is fine. Robot arms painting or welding: PMSM.
\end{notebox}
\end{column}
\end{columns}
\end{frame}

% ---------------------------------------------------------------------
\begin{frame}{Pole pairs: electrical versus mechanical angle}
\begin{columns}[T]
\begin{column}{.42\textwidth}
\centering
\begin{tikzpicture}[scale=.75]
\foreach \cx/\pp/\lab in {0/1/{$p=1$ (2 poles)},3.4/2/{$p=2$ (4 poles)}}{
\begin{scope}[shift={(\cx,0)}]
\pgfmathsetmacro{\seg}{180/\pp}
\foreach \k in {1,...,\pp}{
  \pgfmathsetmacro{\aa}{(2*\k-2)*\seg} \pgfmathsetmacro{\ab}{(2*\k-1)*\seg} \pgfmathsetmacro{\ac}{(2*\k)*\seg}
  \fill[mkRed!70] (0,0) -- (\aa:1.3) arc (\aa:\ab:1.3) -- cycle;
  \fill[mkMid] (0,0) -- (\ab:1.3) arc (\ab:\ac:1.3) -- cycle;}
\fill[mkGray] (0,0) circle (0.2);
\node[font=\scriptsize,text=mkNavy] at (0,-1.7) {\lab};
\end{scope}}
\end{tikzpicture}

\vspace{1mm}
\small
One N--S pair passing a coil $=$ one \hl{electrical} cycle ($360^\circ_\text{el}$).
\end{column}
\begin{column}{.55\textwidth}
\small
\begin{eqbox}
$\theta_\text{el}=p\,\theta_\text{mech}$, \qquad $\omega_\text{el}=p\,\omega_\text{mech}$, \qquad $f_\text{el}=\dfrac{p\,n}{60}$
\end{eqbox}
\begin{defbox}[Worked example: robot joint motor]
\small
\begin{enumerate}\setlength\itemsep{1pt}
\item PMSM with $p=4$ pole pairs (8 poles) at $n=\SI{3000}{rpm}$.
\item Mechanical: $3000/60=\SI{50}{rev\per\second}$.
\item Electrical: $f_\text{el}=4\cdot50=\SI{200}{\hertz}$.
\item The drive must make \SI{200}{\hertz} sine currents.
\end{enumerate}
\end{defbox}
\begin{itemize}\setlength\itemsep{1pt}
\item More poles $\Rightarrow$ more torque from a flat, short motor (good for robot joints).
\item But higher electrical frequency $\Rightarrow$ faster drive needed.
\end{itemize}
\end{column}
\end{columns}
\end{frame}

% ---------------------------------------------------------------------
\begin{frame}{Inside a robot joint: frameless motor, brake, encoder, gear}
\begin{columns}[T]
\begin{column}{.6\textwidth}
\centering
\begin{tikzpicture}[scale=.85,lbl/.style={font=\scriptsize,text=mkNavy,align=center},ld/.style={mkGray,line width=.3pt}]
\draw[mkGray,line width=1pt,fill=mkLight] (0,0) rectangle (8,3);
\fill[white] (0.2,1.2) rectangle (7.8,1.8);
\node[font=\tiny,text=mkGray] at (4,1.5) {hollow shaft: cables pass through};
% encoder
\fill[mkGreen!60] (0.4,1.85) rectangle (0.7,2.8); \fill[mkGreen!60] (0.4,0.2) rectangle (0.7,1.15);
% brake
\fill[mkAmber!70] (1.0,1.85) rectangle (1.5,2.8); \fill[mkAmber!70] (1.0,0.2) rectangle (1.5,1.15);
% motor stator + rotor
\fill[mkSky!60] (1.9,2.2) rectangle (4.3,2.85); \fill[mkSky!60] (1.9,0.15) rectangle (4.3,0.8);
\fill[mkRed!60] (1.9,1.85) rectangle (4.3,2.1); \fill[mkMid] (1.9,0.9) rectangle (4.3,1.15);
% gear
\fill[mkOrange!60] (4.8,1.85) rectangle (6.4,2.85); \fill[mkOrange!60] (4.8,0.15) rectangle (6.4,1.15);
\fill[mkGray!70] (6.4,1.85) rectangle (7.9,2.4); \fill[mkGray!70] (6.4,0.6) rectangle (7.9,1.15);
\draw[-{Stealth},mkNavy,line width=1pt] (8.3,2.5) arc (90:-200:0.35);
\draw[ld] (0.55,2.8) -- (0.4,3.6) node[above,lbl]{encoder};
\draw[ld] (1.25,2.8) -- (1.6,3.6) node[above,lbl]{brake};
\draw[ld] (3.1,2.85) -- (3.1,3.6) node[above,lbl]{frameless PMSM\\ (stator + rotor)};
\draw[ld] (5.6,2.85) -- (5.6,3.6) node[above,lbl]{gearbox\\ (e.g.\ harmonic)};
\draw[ld] (7.2,0.6) -- (7.2,-0.4) node[below,lbl]{output flange $\rightarrow$ next link};
\end{tikzpicture}
\end{column}
\begin{column}{.37\textwidth}
\small
\begin{itemize}\setlength\itemsep{2pt}
\item \term{Frameless motor}: only stator ring and rotor ring are bought; the robot housing is the motor housing.
\item \term{Hollow shaft}: cables run through the joint.
\item Order: encoder $\rightarrow$ brake $\rightarrow$ motor $\rightarrow$ gearbox $\rightarrow$ link.
\item Compact, light, integrated: typical of \hl{cobots}.
\end{itemize}
\gofurther{many cobot joints add a second encoder on the output side to measure the gearbox twist (W2D4 sensors).}
\end{column}
\end{columns}
\end{frame}

% ---------------------------------------------------------------------
\begin{frame}{Comparison: DC, BLDC and PMSM in robots}
\centering
\small
\renewcommand{\arraystretch}{1.15}
\begin{tabularx}{\textwidth}{@{}l*{3}{>{\centering\arraybackslash}X}@{}}
\toprule
 & \textbf{Brushed DC} & \textbf{BLDC} & \textbf{PMSM (AC servo)}\\
\midrule
Commutation & brushes + commutator & 6 steps, Hall sensors & sine currents, encoder\\
Coils on & rotor & stator & stator\\
Torque per size & medium & high & \textcolor{mkGreen}{\textbf{very high}}\\
Smoothness & medium & \textcolor{mkRed}{ripple at each step} & \textcolor{mkGreen}{\textbf{very smooth}}\\
Maintenance & \textcolor{mkRed}{brush wear} & none & none\\
Drive electronics & \textcolor{mkGreen}{\textbf{simplest}} & simple & most complex\\
Robot examples & toy robots, small grippers & drones, AMR wheels, fans & arm joints, cobots, CNC axes\\
\bottomrule
\end{tabularx}

\vspace{4mm}
\begin{keybox}[Rule of thumb]
\small
Need precise, smooth, dynamic motion? $\Rightarrow$ \hl{PMSM}. Need cheap and simple? $\Rightarrow$ brushed DC or BLDC.
BLDC and PMSM are built almost the same -- the difference is mainly \hl{how the drive feeds them}.
\end{keybox}
\end{frame}

% ---------------------------------------------------------------------
\begin{frame}{Session 1.2 -- Check your understanding}
\begin{columns}[T]
\begin{column}{.62\textwidth}
\begin{keybox}[Check your understanding]
\begin{enumerate}\setlength\itemsep{6pt}
\item What job does the commutator do in a brushed DC motor -- and what does that job in a BLDC motor?
\item Why is a brushless motor easier to cool than a brushed motor of the same size?
\item A PMSM with $p=5$ pole pairs runs at \SI{1200}{rpm}. What electrical frequency must the drive produce?
\end{enumerate}
\end{keybox}
\end{column}
\begin{column}{.34\textwidth}
\small
\textbf{Hint:} think about \emph{where} the current flows (rotor or stator), \emph{who} knows the rotor angle, and $f_\text{el}=p\,n/60$.
\end{column}
\end{columns}
\end{frame}

% ---------------------------------------------------------------------
\begin{frame}{Session 1.2 -- Summary and key terms}
\begin{columns}[T]
\begin{column}{.55\textwidth}
\small
\begin{itemize}
\item Every motor: stator, rotor, air gap; torque is largest with field and magnet $90^\circ$ apart.
\item Coil currents must be switched as the rotor turns (\term{commutation}).
\item Brushed DC: mechanical switch. BLDC: 6 steps + Hall. PMSM: sine + encoder.
\item More pole pairs: more torque, higher electrical frequency.
\end{itemize}
\begin{keybox}[Remember]
$\omega_\text{el}=p\,\omega_\text{mech}$, \quad $f_\text{el}=p\,n/60$
\end{keybox}
\end{column}
\begin{column}{.42\textwidth}
\begin{defbox}[Key terms]
\small
\term{Stator / rotor}: fixed / turning part.\\
\term{Commutation}: switching coil currents with rotor angle.\\
\term{Hall sensor}: detects magnet polarity (on/off).\\
\term{Torque ripple}: periodic torque wobble.\\
\term{Pole pair $p$}: one N + one S magnet.\\
\term{Frameless motor}: stator + rotor without housing.
\end{defbox}
\end{column}
\end{columns}
\end{frame}

% =====================================================================
\section{1.3 Electrical and mechanical motor models}
% =====================================================================

\begin{frame}{Refresher: Newton's law for rotation}
\begin{columns}[T]
\begin{column}{.42\textwidth}
\centering
\begin{tikzpicture}[scale=.9]
\fill[mkSky!60,draw=mkBlue] (0,0) circle (1.3);
\fill[mkGray] (0,0) circle (0.15);
\draw[-{Stealth},mkOrange,line width=1.4pt] (1.6,-0.6) arc (-20:60:1.7) node[above,font=\small,text=mkOrange]{$T$ (push)};
\draw[-{Stealth},mkNavy,line width=1pt] (-1.55,0.6) arc (160:230:1.6) node[below left,font=\small]{$\omega$};
\node[font=\small] at (0,-1.7) {disc with inertia $J$};
\end{tikzpicture}

\vspace{1mm}
\small
\renewcommand{\arraystretch}{1.1}
\begin{tabular}{@{}lll@{}}
\toprule
\textbf{Linear} & \textbf{Rotation} & \textbf{Unit}\\
\midrule
mass $m$ & inertia $J$ & \si{\kilo\gram\metre\squared}\\
$F=m\,\dfrac{dv}{dt}$ & $T=J\,\dfrac{d\omega}{dt}$ & \\
\bottomrule
\end{tabular}
\end{column}
\begin{column}{.55\textwidth}
\small
\textbf{Analogy:} a heavy playground roundabout is hard to start and hard to stop -- it has a large \term{moment of inertia} $J$.
\begin{eqbox}
$J\,\dfrac{d\omega}{dt} = \sum T = T_\text{motor} - T_\text{friction} - T_\text{load}$
\end{eqbox}
\begin{defbox}[Worked example]
\small
\begin{enumerate}\setlength\itemsep{1pt}
\item Rotor $J=\SI{1e-4}{\kilo\gram\metre\squared}$, net torque $T=\SI{0.5}{\newton\metre}$.
\item Acceleration $\alpha=T/J=\SI{5000}{\radian\per\second\squared}$.
\item Time to reach \SI{314}{\radian\per\second} (3000 rpm): $t=314/5000=\SI{0.063}{\second}$.
\end{enumerate}
\end{defbox}
Disc: $J=\tfrac12 m r^2$. \hl{Take-away:} net torque sets how fast the speed changes.
\end{column}
\end{columns}
\end{frame}

% ---------------------------------------------------------------------
\begin{frame}{Refresher: first-order ODEs and the exponential response}
\begin{columns}[T]
\begin{column}{.5\textwidth}
\begin{tikzpicture}
\begin{axis}[mkplot,width=7.2cm,height=4.8cm,xmin=0,xmax=5,ymin=0,ymax=1.1,
  xlabel={time $t/\tau$},ylabel={$y/y_\infty$},xtick={0,1,2,3,4,5}]
\addplot[mkBlue,thick,domain=0:5,samples=80]{1-exp(-x)};
\addplot[mkOrange,thick,dashed,domain=0:1]{x};
\draw[dashed,mkGray] (axis cs:1,0) -- (axis cs:1,0.632) -- (axis cs:0,0.632);
\draw[dashed,mkGray] (axis cs:3,0) -- (axis cs:3,0.95);
\node[font=\scriptsize,anchor=west] at (axis cs:1.05,0.55) {63\,\% after $1\tau$};
\node[font=\scriptsize,anchor=west] at (axis cs:3.05,0.85) {95\,\% after $3\tau$};
\node[font=\scriptsize,text=mkOrange,anchor=west] at (axis cs:0.15,0.98) {initial slope};
\end{axis}
\end{tikzpicture}
\end{column}
\begin{column}{.48\textwidth}
\small
\textbf{You know this from RC charging:} the capacitor voltage rises fast at first, then slower and slower.
\begin{eqbox}
$\tau\dfrac{dy}{dt}+y=y_\infty \;\Rightarrow\; y(t)=y_\infty\big(1-e^{-t/\tau}\big)$
\end{eqbox}
\begin{itemize}\setlength\itemsep{1pt}
\item $\tau$ = \term{time constant} [s]: RC circuit $\tau=RC$, RL circuit $\tau=L/R$.
\item After $1\tau$: 63\,\%; after $3\tau$: 95\,\%; after $5\tau$: ``done''.
\end{itemize}
\textbf{Example:} $L=\SI{0.5}{\milli\henry}$, $R=\SI{1}{\ohm}$ $\Rightarrow$ $\tau=\SI{0.5}{\milli\second}$; current settles in $\approx\SI{1.5}{\milli\second}$.\\[1mm]
\hl{Take-away:} one energy store $\Rightarrow$ one $\tau$ $\Rightarrow$ one exponential.
\end{column}
\end{columns}
\end{frame}

% ---------------------------------------------------------------------
\begin{frame}{Equivalent circuit of the DC motor}
\begin{columns}[T]
\begin{column}{.55\textwidth}
\centering
\begin{circuitikz}[scale=.64,transform shape,font=\small]
\draw (0,0) to[V,v<=$v$,color=mkNavy] (0,3)
  to[R,l=$R$,i>^=$i$,color=mkNavy] (2.6,3)
  to[L,l=$L$,color=mkNavy] (4.8,3)
  to[short] (5.6,3)
  to[sV,v^<={$e=k_e\omega$},color=mkOrange] (5.6,0) -- (0,0);
\draw[mkGray,dashed] (7.0,-0.4) -- (7.0,3.4);
\begin{scope}[shift={(8.3,1.5)}]
\fill[mkSky!60] (0,-0.9) rectangle (0.5,0.9);
\node[font=\scriptsize] at (0.25,-1.15) {$J$};
\draw[mkGray,line width=2pt] (0.5,0) -- (1.9,0);
\draw[mkGray,line width=2pt] (-0.7,0) -- (0,0);
\draw[-{Stealth},mkOrange,line width=.9pt] (-0.4,0.6) arc (120:240:0.65);
\node[font=\scriptsize,text=mkOrange] at (-0.7,0.9) {$T=k_t i$};
\draw[-{Stealth},mkNavy,line width=.9pt] (1.5,-0.6) arc (-120:120:0.65);
\node[font=\scriptsize] at (2.0,0.95) {$T_L$};
\node[font=\scriptsize] at (1.2,-1.0) {$\omega,\;B\omega$};
\end{scope}
\node[font=\scriptsize,text=mkGray] at (2.8,-0.8) {electrical side};
\node[font=\scriptsize,text=mkGray] at (9.0,-0.8) {mechanical side};
\end{circuitikz}

\vspace{2mm}
\small
The electrical side is the \hl{RL circuit} you know, plus one extra ``battery'' $e=k_e\omega$ that grows with speed and \hl{opposes} the supply.
\end{column}
\begin{column}{.43\textwidth}
\small
\begin{tabular}{@{}lll@{}}
\toprule
\textbf{Sym.} & \textbf{Meaning} & \textbf{Unit}\\
\midrule
$R$ & winding resistance & \si{\ohm}\\
$L$ & winding inductance & \si{\henry}\\
$k_e$ & back-EMF constant & \si{\volt\second\per\radian}\\
$k_t$ & torque constant & \si{\newton\metre\per\ampere}\\
$J$ & rotor (+ load) inertia & \si{\kilo\gram\metre\squared}\\
$B$ & viscous friction & \si{\newton\metre\second\per\radian}\\
$T_L$ & load torque & \si{\newton\metre}\\
\bottomrule
\end{tabular}

\vspace{2mm}
{\small\textbf{Signs:} $v,i$ positive into the motor; $T_L$ works against the motion.}\\[1mm]
{\small BLDC/PMSM: same model per phase (W1D5).}
\end{column}
\end{columns}
\end{frame}

% ---------------------------------------------------------------------
\begin{frame}{The two motor equations}
\begin{columns}[T]
\begin{column}{.52\textwidth}
\small
\textbf{Electrical (KVL around the loop):}
\[
v = R\,i + L\frac{di}{dt} + k_e\,\omega
\]
\textbf{Mechanical (Newton for rotation):}
\begin{eqbox}
$J\dfrac{d\omega}{dt} = k_t\,i - B\,\omega - T_L$
\end{eqbox}
\begin{itemize}\setlength\itemsep{1pt}
\item Linked both ways: $i$ makes torque, $\omega$ makes back-EMF.
\item Two first-order ODEs $\Rightarrow$ one \hl{second-order} system.
\end{itemize}
\end{column}
\begin{column}{.45\textwidth}
\begin{keybox}[Why the two constants are equal]
\small
Power taken from the circuit by the back-EMF: $e\,i=k_e\,\omega\,i$\\
Power delivered to the shaft: $T\,\omega=k_t\,i\,\omega$\\
Energy is conserved $\Rightarrow$ $\bm{k_t=k_e}$ (SI units).
\end{keybox}
\small
\textbf{Datasheet units:} $k_e$ is often given in V per 1000 rpm.
Divide by 104.7 to get \si{\volt\second\per\radian}:
\SI{10.47}{\volt}/krpm $\Rightarrow$ $k_e=k_t=0.1$ SI.\\[1mm]
\gofurther{three-phase datasheets use rms and line-to-line values, so $k_t$ and $k_e$ differ by factors like $\sqrt3$; resistance rises $\approx$0.4\,\%/K when hot.}
\end{column}
\end{columns}
\end{frame}

% ---------------------------------------------------------------------
\begin{frame}{Reading a motor datasheet: worked example}
\begin{columns}[T]
\begin{column}{.45\textwidth}
\small
\renewcommand{\arraystretch}{1.1}
\begin{tabular}{@{}lr@{}}
\toprule
\textbf{Datasheet value (example)} & \\
\midrule
Rated voltage $v_N$ & \SI{48}{\volt}\\
Rated torque $T_N$ & \SI{0.64}{\newton\metre}\\
Rated speed $n_N$ & \SI{3000}{rpm}\\
Torque constant $k_t$ & \SI{0.12}{\newton\metre\per\ampere}\\
Terminal resistance $R$ & \SI{0.8}{\ohm}\\
Terminal inductance $L$ & \SI{1.2}{\milli\henry}\\
Rotor inertia $J_m$ & \SI{0.6e-4}{\kilo\gram\metre\squared}\\
\bottomrule
\end{tabular}\\[1mm]
{\scriptsize\color{mkGray} Illustrative values for a \SI{200}{\watt} class motor.}
\end{column}
\begin{column}{.52\textwidth}
\begin{defbox}[What can we compute?]
\small
\begin{enumerate}\setlength\itemsep{2pt}
\item Rated current: $i_N=T_N/k_t=0.64/0.12\approx\SI{5.3}{\ampere}$.
\item Back-EMF at \SI{3000}{rpm} ($\SI{314}{\radian\per\second}$): $e=0.12\cdot314\approx\SI{38}{\volt}$.
\item Voltage needed: $v=Ri+e=0.8\cdot5.3+38\approx\SI{42}{\volt}<\SI{48}{\volt}$ \checkmark
\item Copper loss: $P_\text{Cu}=Ri^2=0.8\cdot5.3^2\approx\SI{23}{\watt}$ of heat.
\end{enumerate}
\end{defbox}
\small
Output power: $T\omega=0.64\cdot314\approx\SI{200}{\watt}$.\\
\hl{Take-away:} back-EMF uses most of the supply voltage at full speed.
\end{column}
\end{columns}
\end{frame}

% ---------------------------------------------------------------------
\begin{frame}{Friction and load torque}
\begin{columns}[T]
\begin{column}{.48\textwidth}
\begin{tikzpicture}
\begin{axis}[mkplot,width=7cm,height=5.2cm,xmin=-3,xmax=3,ymin=-1.9,ymax=1.9,
  axis lines=middle,xlabel={$\omega$},ylabel={$T_f$},xtick=\empty,ytick=\empty,
  every axis x label/.style={at={(current axis.right of origin)},anchor=west,font=\scriptsize},
  every axis y label/.style={at={(current axis.above origin)},anchor=south,font=\scriptsize},
  legend style={at={(0.02,0.98)},anchor=north west}]
\addplot[mkSky,thick,dashed,domain=0.001:3,samples=40]{0.25*x}; \addlegendentry{viscous $B\omega$}
\addplot[mkMid,thick,dotted,domain=0.001:3,samples=2]{0.8}; \addlegendentry{Coulomb (constant)}
\addplot[mkOrange,thick,domain=0.001:3,samples=120]{0.8+0.6*exp(-(x/0.35)^2)+0.25*x}; \addlegendentry{real friction}
\addplot[mkSky,thick,dashed,domain=-3:-0.001,samples=40,forget plot]{0.25*x};
\addplot[mkMid,thick,dotted,domain=-3:-0.001,samples=2,forget plot]{-0.8};
\addplot[mkOrange,thick,domain=-3:-0.001,samples=120,forget plot]{-0.8-0.6*exp(-(x/0.35)^2)+0.25*x};
\draw[mkOrange,thick] (axis cs:0,-1.4) -- (axis cs:0,1.4);
\node[font=\scriptsize,anchor=west,text=mkOrange] at (axis cs:0.05,1.48) {stiction};
\end{axis}
\end{tikzpicture}
\end{column}
\begin{column}{.5\textwidth}
\small
\textbf{Analogy:} pushing a heavy box -- hard to get it moving, easier once it slides.
\begin{itemize}\setlength\itemsep{2pt}
\item \term{Viscous} friction $B\omega$: grows with speed, like stirring honey.
\item \term{Coulomb} friction: constant, flips sign with direction.
\item \term{Stiction}: extra ``sticking'' torque at standstill $\Rightarrow$ jerky start.
\item \term{Load torque} $T_L$: gravity on a robot arm, payload, contact forces.
\end{itemize}
\begin{eqbox}
model used today: $T_f\approx B\,\omega$ \quad (linear, easy maths)
\end{eqbox}
{\scriptsize\color{mkGray}[Armstrong-H\'elouvry, Dupont \& Canudas de Wit, 1994]}
\end{column}
\end{columns}
\end{frame}

% ---------------------------------------------------------------------
\begin{frame}{Steady state: the speed--torque line}
\begin{columns}[T]
\begin{column}{.58\textwidth}
\begin{tikzpicture}
\begin{axis}[mkplot,width=7.6cm,height=5.0cm,xmin=0,xmax=5,ymin=0,ymax=5000,
  xlabel={load torque $T$ (\si{\newton\metre})},ylabel={speed $n$ (rpm)},scaled y ticks=false,
  yticklabel style={/pgf/number format/fixed},legend pos=north east]
\addplot[mkBlue,thick,domain=0:4.8]{(480-100*x)*9.549}; \addlegendentry{$v=\SI{48}{\volt}$}
\addplot[mkSky,thick,dashed,domain=0:2.4]{(240-100*x)*9.549}; \addlegendentry{$v=\SI{24}{\volt}$}
\node[circle,fill=mkOrange,inner sep=1.4pt] at (axis cs:0,4584) {};
\node[font=\scriptsize,anchor=west,text=mkOrange] at (axis cs:0.1,4700) {no-load speed};
\node[circle,fill=mkOrange,inner sep=1.4pt] at (axis cs:4.8,0) {};
\node[font=\scriptsize,anchor=south east,text=mkOrange] at (axis cs:4.8,150) {stall};
\end{axis}
\end{tikzpicture}\\
{\scriptsize\color{mkGray} Example motor: $R=\SI{1}{\ohm}$, $k_t=k_e=0.1$ SI, friction neglected.}
\end{column}
\begin{column}{.4\textwidth}
\small
Constant speed: set $\tfrac{d}{dt}=0$ in both equations:
\begin{eqbox}
$\omega=\dfrac{v}{k_e}-\dfrac{R}{k_e k_t}\,T$
\end{eqbox}
\begin{enumerate}\setlength\itemsep{1pt}
\item No load: $\omega_0=48/0.1=\SI{480}{\radian\per\second}$ (\SI{4584}{rpm}).
\item Stalled: $i=v/R=\SI{48}{\ampere}$, $T=k_t i=\SI{4.8}{\newton\metre}$.
\item More load $\Rightarrow$ speed \hl{droops} along the line.
\item Higher voltage $\Rightarrow$ line moves up.
\end{enumerate}
\end{column}
\end{columns}
\end{frame}

% ---------------------------------------------------------------------
\begin{frame}{Output power and efficiency along the line}
\begin{columns}[T]
\begin{column}{.56\textwidth}
\begin{tikzpicture}
\begin{axis}[mkplot,width=7.6cm,height=5.0cm,xmin=0,xmax=4.8,ymin=0,ymax=650,
  xlabel={load torque $T$ (\si{\newton\metre})},ylabel={$P_\text{out}$ (W)},axis y line*=left,axis x line*=bottom]
\addplot[mkBlue,thick,domain=0:4.7,samples=100]{x*(470-100*x)}; \label{pl:pout}
\draw[dashed,mkGray] (axis cs:2.35,0) -- (axis cs:2.35,552);
\node[font=\scriptsize,text=mkGray,anchor=south] at (axis cs:2.35,560) {max power at $\approx\tfrac12$ stall};
\end{axis}
\begin{axis}[mkplot,width=7.6cm,height=5.0cm,xmin=0,xmax=4.8,ymin=0,ymax=100,grid=none,
  axis y line*=right,axis x line=none,ylabel={$\eta$ (\%)},legend pos=north east,legend style={at={(0.98,0.80)}}]
\addlegendimage{/pgfplots/refstyle=pl:pout}\addlegendentry{output power}
\addplot[mkOrange,thick,domain=0:4.7,samples=150]{100*x*(470-100*x)/(480*(x+0.1))}; \addlegendentry{efficiency}
\end{axis}
\end{tikzpicture}\\
{\scriptsize\color{mkGray} Same motor at \SI{48}{\volt}, with \SI{0.1}{\newton\metre} friction torque.}
\end{column}
\begin{column}{.42\textwidth}
\small
\begin{itemize}\setlength\itemsep{2pt}
\item $P_\text{out}=T\,\omega$ is a \hl{parabola}: zero at no load (no torque) and at stall (no speed).
\item Efficiency is \hl{highest at light load}; it falls as copper loss $Ri^2$ grows.
\item Real motors are rated far below stall -- heat (not the line) sets the limit.
\end{itemize}
\begin{eqbox}
$\eta=\dfrac{T\,\omega}{v\,i}$, \quad losses $=v\,i-T\omega$ $\rightarrow$ heat
\end{eqbox}
\gofurther{rated (continuous) torque is set by winding temperature; peak torque by magnet demagnetisation and drive current (W1D2 sizing).}
\end{column}
\end{columns}
\end{frame}

% ---------------------------------------------------------------------
\begin{frame}{Block diagram of the DC motor}
\centering
\begin{tikzpicture}[node distance=7mm]
\node (v) {$v$};
\node[mksum,right=6mm of v] (s1) {};
\node[mkblk,right=of s1,minimum width=18mm] (el) {electrical\\ RL circuit};
\node[mkblk,right=of el] (kt) {$\times k_t$};
\node[mksum,right=of kt] (s2) {};
\node[mkblk,right=of s2,minimum width=18mm] (me) {inertia $J$\\ + friction $B$};
\node[mkblk,right=10mm of me] (int) {integrate\\ $\int dt$};
\node[right=6mm of int] (th) {$\theta$};
\node[above=6mm of s2] (tl) {$T_L$};
\node[mkblk,below=9mm of kt,fill=mkCream,draw=mkOrange] (ke) {$\times k_e$};
\draw[mksig] (v) -- node[pos=.85,above,font=\scriptsize]{$+$} (s1);
\draw[mksig] (s1) -- (el);
\draw[mksig] (el) -- node[above,mklab]{$i$} (kt);
\draw[mksig] (kt) -- node[above,mklab]{$T$} node[pos=.85,above,font=\scriptsize]{$+$} (s2);
\draw[mksig] (tl) -- node[pos=.8,right,font=\scriptsize]{$-$} (s2);
\draw[mksig] (s2) -- (me);
\draw[mksig] (me) -- node[above,mklab]{$\omega$} coordinate[pos=.4] (wtap) (int);
\draw[mksig] (int) -- (th);
\draw[mksig] (wtap) |- (ke);
\draw[mksig] (ke) -| node[pos=.9,left,font=\scriptsize]{$-$} node[pos=.3,below,mklab]{$e$} (s1);
\end{tikzpicture}

\vspace{4mm}
\begin{columns}[T]
\begin{column}{.5\textwidth}
\small
\begin{itemize}\setlength\itemsep{2pt}
\item A \term{block diagram} draws the equations as signal flow: boxes act on signals, circles add or subtract.
\item Back-EMF is a built-in \term{feedback loop}: faster $\Rightarrow$ more $e$ $\Rightarrow$ less current $\Rightarrow$ less torque.
\item $T_L$ (gravity, payload) enters as a \hl{disturbance}.
\end{itemize}
\end{column}
\begin{column}{.46\textwidth}
\begin{notebox}
\small
Robot drives control the \textbf{current} directly, so $T\approx k_t\,i^*$: the motor behaves like a \hl{torque source} (W1D4--W1D5).
\end{notebox}
\gofurther{with the Laplace transform each box becomes a transfer function, e.g.\ $1/(Ls+R)$ and $1/(Js+B)$ (W1D2).}
\end{column}
\end{columns}
\end{frame}

% ---------------------------------------------------------------------
\begin{frame}{The motor's two time constants: worked example}
\begin{columns}[T]
\begin{column}{.55\textwidth}
\begin{tikzpicture}
\begin{axis}[mkplot,width=7.6cm,height=5.0cm,xmin=0,xmax=50,ymin=0,ymax=260,
  xlabel={$t$ (ms)},ylabel={$\omega$ (\si{\radian\per\second})},axis y line*=left,axis x line*=bottom]
\addplot[mkBlue,thick,domain=0:50,samples=200]{238.81+14.169*exp(-1.8944*x)-252.97*exp(-0.10610*x)}; \label{pl:w}
\draw[dashed,mkGray] (axis cs:9.42,0) -- (axis cs:9.42,151) -- (axis cs:0,151);
\node[font=\scriptsize,text=mkGray,anchor=west] at (axis cs:9.8,140) {63\,\% at $\approx\tau_m$};
\end{axis}
\begin{axis}[mkplot,width=7.6cm,height=5.0cm,xmin=0,xmax=50,ymin=0,ymax=26,grid=none,
  axis y line*=right,axis x line=none,ylabel={$i$ (A)},legend pos=north east]
\addlegendimage{/pgfplots/refstyle=pl:w}\addlegendentry{speed $\omega$}
\addplot[mkOrange,thick,domain=0:50,samples=300]{0.1194-26.834*exp(-1.8944*x)+26.715*exp(-0.10610*x)}; \addlegendentry{current $i$}
\end{axis}
\end{tikzpicture}\\
{\scriptsize\color{mkGray} Simulated \SI{24}{\volt} step, $R=\SI{1}{\ohm}$, $L=\SI{0.5}{\milli\henry}$, $k_t=k_e=0.1$, $J=\SI{1e-4}{\kilo\gram\metre\squared}$.}
\end{column}
\begin{column}{.43\textwidth}
\small
\begin{enumerate}\setlength\itemsep{2pt}
\item Electrical: $\tau_e=L/R=\SI{0.5}{\milli\second}$.
\item Mechanical: $\tau_m=\dfrac{RJ}{k_tk_e}=\dfrac{1\cdot10^{-4}}{0.01}=\SI{10}{\milli\second}$.
\item Final speed: $\omega_\infty\approx v/k_e=\SI{240}{\radian\per\second}$.
\item Start current: $\omega=0$, no back-EMF $\Rightarrow$ up to $v/R=\SI{24}{\ampere}$ (peak \SI{21}{\ampere}).
\end{enumerate}
\begin{eqbox}
$\tau_e=\dfrac{L}{R}$ \quad $\ll$ \quad $\tau_m=\dfrac{RJ}{k_tk_e}$
\end{eqbox}
\hl{Take-away:} current is fast, speed is slow. Real drives limit the start current.
\end{column}
\end{columns}
\end{frame}

% ---------------------------------------------------------------------
\begin{frame}{What the robot arm adds: load inertia slows the motor}
\begin{columns}[T]
\begin{column}{.55\textwidth}
\begin{tikzpicture}
\begin{axis}[mkplot,width=7.6cm,height=5.0cm,xmin=0,xmax=150,ymin=0,ymax=260,
  xlabel={$t$ (ms)},ylabel={$\omega$ (\si{\radian\per\second})},legend pos=south east]
\addplot[mkBlue,thick,domain=0:150,samples=150]{240*(1-exp(-x/10))}; \addlegendentry{motor alone, $\tau_m=\SI{10}{\milli\second}$}
\addplot[mkSky,thick,dashed,domain=0:150,samples=150]{240*(1-exp(-x/20))}; \addlegendentry{$J$ doubled, \SI{20}{\milli\second}}
\addplot[mkOrange,thick,domain=0:150,samples=150]{240*(1-exp(-x/40))}; \addlegendentry{$J\times4$, \SI{40}{\milli\second}}
\end{axis}
\end{tikzpicture}\\
{\scriptsize\color{mkGray} Same \SI{24}{\volt} step, $\tau_e$ neglected: $\omega(t)=\omega_\infty(1-e^{-t/\tau_m})$.}
\end{column}
\begin{column}{.43\textwidth}
\small
\begin{itemize}\setlength\itemsep{2pt}
\item The motor must accelerate its own rotor \hl{plus} everything attached: gearbox, link, payload.
\item $\tau_m\propto J_\text{total}$: double the inertia, double the time to speed up.
\item The final speed stays the same (set by $v/k_e$).
\item A robot arm stretched out has more inertia than when folded $\Rightarrow$ \hl{its behaviour changes with pose}.
\end{itemize}
\begin{notebox}
\small
Gearboxes shrink the load inertia the motor ``feels'' by $N^2$ -- the topic of \term{reflected inertia} (W1D2).
\end{notebox}
\end{column}
\end{columns}
\end{frame}

% ---------------------------------------------------------------------
\begin{frame}{Worked example: holding a robot arm against gravity}
\begin{columns}[T]
\begin{column}{.4\textwidth}
\centering
\begin{tikzpicture}[scale=.95]
\fill[mkGray!50] (-0.6,-0.6) rectangle (0.6,-0.3);
\fill[mkSky!60,draw=mkBlue] (0,0) circle (0.35);
\draw[mkBlue,line width=5pt] (0,0) -- (3.2,0);
\fill[mkOrange] (3.2,0) circle (0.25);
\draw[-{Stealth},mkRed,line width=1.2pt] (3.2,-0.25) -- (3.2,-1.4) node[below,font=\scriptsize]{$m g$};
\draw[|-|,mkNavy] (0,0.55) -- node[above,font=\scriptsize]{$r=\SI{0.5}{\metre}$} (3.2,0.55);
\node[font=\scriptsize,text=mkNavy,align=center] at (0,-1.1) {joint:\\ motor + gear};
\node[font=\scriptsize,text=mkOrange] at (3.2,0.3) {\phantom{x}};
\end{tikzpicture}

\vspace{2mm}
\small
Payload $m=\SI{2}{\kilo\gram}$ at the end of a horizontal link (link mass neglected).
\end{column}
\begin{column}{.57\textwidth}
\begin{defbox}[Step by step]
\small
\begin{enumerate}\setlength\itemsep{2pt}
\item Joint torque: $T_L=m g r=2\cdot9.81\cdot0.5\approx\SI{9.8}{\newton\metre}$.
\item Direct drive, $k_t=\SI{0.1}{\newton\metre\per\ampere}$: $i=9.8/0.1=$ \hl{\SI{98}{\ampere}} -- far too much!
\item With a 100:1 gearbox the motor sees $9.8/100=\SI{0.098}{\newton\metre}$.
\item Current: $i=0.098/0.1\approx\SI{1}{\ampere}$; copper loss $Ri^2\approx\SI{1}{\watt}$ ($R=\SI{1}{\ohm}$).
\end{enumerate}
\end{defbox}
\small
\begin{itemize}\setlength\itemsep{1pt}
\item Holding still needs \hl{torque, not power}: $\omega=0$ so $T\omega=0$, yet current flows and heats the motor.
\item That is why robot joints have gearboxes and a \hl{holding brake}.
\end{itemize}
\end{column}
\end{columns}
\end{frame}

% ---------------------------------------------------------------------
\begin{frame}{Session 1.3 -- Check your understanding}
\begin{columns}[T]
\begin{column}{.62\textwidth}
\begin{keybox}[Check your understanding]
\begin{enumerate}\setlength\itemsep{6pt}
\item A motor has $\tau_m=\SI{20}{\milli\second}$. Roughly how long until it reaches 95\,\% of its final speed after a voltage step?
\item You attach a load that doubles the total inertia $J$. What happens to $\tau_m$ and to the final speed?
\item Why is the current so large in the first milliseconds after switching on, and why does it then fall?
\end{enumerate}
\end{keybox}
\end{column}
\begin{column}{.34\textwidth}
\small
\textbf{Hint:} use the 63\,\% / 95\,\% rule, the formula $\tau_m=RJ/(k_tk_e)$ and the back-EMF $e=k_e\omega$.
\end{column}
\end{columns}
\end{frame}

% ---------------------------------------------------------------------
\begin{frame}{Session 1.3 -- Summary and key terms}
\begin{columns}[T]
\begin{column}{.55\textwidth}
\small
\begin{itemize}
\item Motor = RL circuit + back-EMF, coupled to an inertia with friction and load.
\item Energy conservation gives $k_t=k_e$ (SI units).
\item Constant speed: straight speed--torque line; power peaks at half stall.
\item Two time constants: fast electrical $\tau_e$, slow mechanical $\tau_m$.
\end{itemize}
\begin{keybox}[Remember]
$v=Ri+L\dfrac{di}{dt}+k_e\omega$, \quad $J\dfrac{d\omega}{dt}=k_ti-B\omega-T_L$\\[1mm]
$\tau_e=L/R$, \quad $\tau_m=RJ/(k_tk_e)$
\end{keybox}
\end{column}
\begin{column}{.42\textwidth}
\begin{defbox}[Key terms]
\small
\term{Moment of inertia $J$}: resistance to changes of rotation speed.\\
\term{Time constant $\tau$}: time to reach 63\,\%.\\
\term{Stiction}: sticking friction at standstill.\\
\term{Stall torque}: torque at zero speed.\\
\term{Block diagram}: equations drawn as signal flow.\\
\term{Disturbance}: unwanted input such as $T_L$.
\end{defbox}
\end{column}
\end{columns}
\end{frame}

% =====================================================================
\section{1.4 Stepper and induction motors in automation}
% =====================================================================

\begin{frame}{Stepper motor: moving in exact steps}
\begin{columns}[T]
\begin{column}{.55\textwidth}
\centering
\begin{tikzpicture}[scale=.55]
\foreach \cx/\ang/\lab/\on in {0/90/{1: A$+$}/90,3.6/0/{2: B$+$}/0,7.2/-90/{3: A$-$}/-90,10.8/180/{4: B$-$}/180}{
\begin{scope}[shift={(\cx,0)}]
\draw[mkGray,line width=2.5pt] (0,0) circle (1.45);
\foreach \a in {0,90,180,270}{ \fill[mkGray!50,rotate=\a] (0.95,-0.22) rectangle (1.4,0.22);}
\fill[mkOrange,rotate=\on] (0.95,-0.22) rectangle (1.4,0.22);
\fill[mkOrange,rotate=\on+180] (0.95,-0.22) rectangle (1.4,0.22);
\begin{scope}[rotate=\ang]
\fill[mkRed!70] (0,-0.17) rectangle (0.8,0.17);
\fill[mkMid] (-0.8,-0.17) rectangle (0,0.17);
\end{scope}
\fill[mkGray] (0,0) circle (0.08);
\node[font=\scriptsize,text=mkNavy] at (0,-1.95) {\lab};
\end{scope}}
\end{tikzpicture}\\[1mm]
{\scriptsize\color{mkGray} Two phases A, B; energised poles orange. Each pulse turns the magnet rotor one step ($90^\circ$ here).}

\vspace{2mm}
\figph{width=.55\linewidth,height=2.6cm,keepaspectratio}{d1_stepper.jpg}\\
\credit{tony\_duell, CC BY 2.0, via Wikimedia Commons}
\end{column}
\begin{column}{.42\textwidth}
\small
\textbf{Analogy:} a clock's second hand -- it jumps tick by tick, and you know the time by counting ticks.
\begin{itemize}\setlength\itemsep{2pt}
\item Each current pulse moves the rotor \hl{one step}.
\item Real \term{hybrid steppers} have many teeth: \hl{200 steps/rev} ($1.8^\circ$).
\item Position $=$ number of pulses: \hl{no sensor needed} (\term{open loop}).
\end{itemize}
\begin{eqbox}
step angle $=\dfrac{360^\circ}{\text{steps per rev}}$
\end{eqbox}
\end{column}
\end{columns}
\end{frame}

% ---------------------------------------------------------------------
\begin{frame}{Full step, half step and microstepping}
\begin{columns}[T]
\begin{column}{.55\textwidth}
\begin{tikzpicture}
\begin{axis}[mkplot,width=7.6cm,height=4.8cm,xmin=0,xmax=360,ymin=-1.25,ymax=1.25,xtick={0,90,...,360},
  xlabel={electrical angle (deg)},ylabel={phase current},ytick={-1,0,1},legend pos=south west,legend columns=2]
\addplot[mkSky,thick,const plot] coordinates {(0,1) (90,-1) (180,-1) (270,1) (360,1)}; \addlegendentry{A full step}
\addplot[mkBlue,thick,const plot,domain=0:360,samples=33] {cos(x)}; \addlegendentry{A, 1/8 micro}
\end{axis}
\end{tikzpicture}\\
{\scriptsize\color{mkGray} Microstepping: phase currents follow a staircase sine (phase B is the same, shifted $90^\circ$).}
\end{column}
\begin{column}{.42\textwidth}
\small
\begin{itemize}\setlength\itemsep{2pt}
\item \term{Full step}: currents only $+I$ or $-I$; jerky, noisy.
\item \term{Half step}: one phase at a time in between $\Rightarrow$ 400 steps/rev.
\item \term{Microstepping}: currents shaped like sine/cosine in small stairs $\Rightarrow$ smooth, quiet.
\end{itemize}
\begin{defbox}[Worked example: 3D-printer axis]
\small
\begin{enumerate}\setlength\itemsep{1pt}
\item 200 steps/rev $\times$ 16 micro $=3200$ microsteps/rev.
\item Lead screw \SI{8}{\milli\metre}/rev: $8/3200=\SI{2.5}{\micro\metre}$ per microstep.
\item \SI{50}{\milli\metre\per\second} needs $50/8\cdot3200=\SI{20000}{pulses\per\second}$.
\end{enumerate}
\end{defbox}
\end{column}
\end{columns}
\end{frame}

% ---------------------------------------------------------------------
\begin{frame}{Stepper limits: torque falls with speed, steps can be lost}
\begin{columns}[T]
\begin{column}{.55\textwidth}
\begin{tikzpicture}
\begin{axis}[mkplot,width=7.6cm,height=5.0cm,xmin=0,xmax=1500,ymin=0,ymax=0.6,
  xlabel={speed (rpm)},ylabel={torque (\si{\newton\metre})},legend pos=north east,
  scaled x ticks=false,xticklabel style={/pgf/number format/fixed}]
\addplot[mkBlue,thick,domain=0:1500,samples=100]{0.5/(1+(x/450)^1.6)}; \addlegendentry{pull-out torque}
\addplot[mkOrange,thick,dashed,domain=0:1500,samples=2]{0.15}; \addlegendentry{load torque}
\fill[mkRed!25,opacity=.6] (axis cs:765,0) rectangle (axis cs:1500,0.6);
\node[font=\scriptsize,text=mkRed,align=center] at (axis cs:1130,0.30) {steps lost\\ (motor stalls)};
\node[font=\scriptsize,text=mkGray,anchor=west] at (axis cs:30,0.54) {holding torque};
\end{axis}
\end{tikzpicture}\\
{\scriptsize\color{mkGray} Typical shape for a NEMA-17 class motor (illustrative values).}
\end{column}
\begin{column}{.42\textwidth}
\small
\begin{itemize}\setlength\itemsep{2pt}
\item Winding inductance slows current rise: at high step rates the current, and so the torque, \hl{collapses}.
\item If load torque $>$ available torque, the rotor slips a step -- and the controller \hl{never knows}.
\item Keep a \hl{safety margin} (often 30--50\,\%) and ramp the speed up slowly.
\item Holding current flows even at rest $\Rightarrow$ steppers run warm.
\end{itemize}
\begin{notebox}
\small
Fix: add an encoder (\term{closed-loop stepper}) -- or use a servo motor.
\end{notebox}
\end{column}
\end{columns}
\end{frame}

% ---------------------------------------------------------------------
\begin{frame}{Induction motor: currents induced in a cage}
\begin{columns}[T]
\begin{column}{.3\textwidth}
\centering
\figph{width=\linewidth,height=3.4cm,keepaspectratio}{d1_squirrel_cage.jpg}\\
\credit{Zureks, CC BY-SA 3.0, via Wikimedia Commons}\\[1mm]
{\scriptsize Squirrel-cage rotor: aluminium bars shorted by end rings, no magnets.}
\end{column}
\begin{column}{.33\textwidth}
\centering
\begin{tikzpicture}[scale=.8]
\draw[mkGray,line width=3pt] (0,0) circle (1.7);
\fill[mkSky!40] (0,0) circle (1.2);
\foreach \a in {0,30,...,330} \fill[mkGray!70] (\a:1.05) circle (0.09);
\foreach \a in {60,90,120} \fill[mkOrange] (\a:1.05) circle (0.09);
\foreach \a in {240,270,300} \fill[mkMid] (\a:1.05) circle (0.09);
\draw[-{Stealth},mkBlue,line width=1.4pt] (-60:1.95) arc (-60:30:1.95) node[right,font=\scriptsize]{field $n_s$};
\draw[-{Stealth},mkOrange,line width=1.1pt] (-50:0.55) arc (-50:20:0.55);
\node[font=\scriptsize,text=mkOrange] at (0,-0.1) {rotor $n$};
\node[font=\scriptsize,text=mkGray] at (0,-2.3) {field runs ahead; bars feel it move};
\end{tikzpicture}
\end{column}
\begin{column}{.34\textwidth}
\small
\begin{itemize}\setlength\itemsep{1pt}
\item Three-phase stator makes a \hl{rotating field}.
\item The field sweeps past the bars and \hl{induces} currents in them (Faraday) $\Rightarrow$ force (Lorentz).
\item Rotor runs slightly slower: \term{slip} $s$.
\end{itemize}
\begin{eqbox}
$n_s=\dfrac{60f}{p}$, \quad $s=\dfrac{n_s-n}{n_s}$
\end{eqbox}
{\small\textbf{Example:} $f=\SI{50}{\hertz}$, $p=2$: $n_s=\SI{1500}{rpm}$; at \SI{1440}{rpm}, $s=4\,\%$.}
\end{column}
\end{columns}
\end{frame}

% ---------------------------------------------------------------------
\begin{frame}{Torque--speed curve of the induction motor}
\begin{columns}[T]
\begin{column}{.56\textwidth}
\begin{tikzpicture}
\begin{axis}[mkplot,width=7.6cm,height=5.0cm,xmin=0,xmax=1550,ymin=0,ymax=1.1,
  xlabel={speed $n$ (rpm)},ylabel={$T/T_\text{max}$},scaled x ticks=false,
  xticklabel style={/pgf/number format/fixed},legend pos=north west]
\addplot[mkBlue,thick,domain=0:1499,samples=200]{2/(((1500-x)/1500)/0.25+0.25/((1500-x)/1500))}; \addlegendentry{motor torque}
\addplot[mkOrange,thick,dashed,domain=0:1550,samples=2]{0.4}; \addlegendentry{conveyor load}
\node[circle,fill=mkOrange,inner sep=1.4pt] at (axis cs:1424,0.4) {};
\node[font=\scriptsize,anchor=north east,text=mkOrange] at (axis cs:1420,0.37) {operating point};
\node[font=\scriptsize,anchor=south,text=mkBlue] at (axis cs:1125,1.0) {breakdown};
\draw[dashed,mkGray] (axis cs:1500,0) -- (axis cs:1500,1.1);
\node[font=\scriptsize,text=mkGray,anchor=north east] at (axis cs:1495,1.08) {$n_s$};
\node[font=\scriptsize,text=mkGray,anchor=south west] at (axis cs:10,0.48) {starting torque};
\end{axis}
\end{tikzpicture}\\
{\scriptsize\color{mkGray} Simplified (Kloss) curve, breakdown slip 25\,\%.}
\end{column}
\begin{column}{.42\textwidth}
\small
\begin{itemize}\setlength\itemsep{2pt}
\item No slip $\Rightarrow$ no induced current $\Rightarrow$ \hl{no torque} at $n_s$.
\item Near $n_s$ the curve is steep: speed barely changes with load (a few \%).
\item Load too big (past \term{breakdown}) $\Rightarrow$ motor stalls.
\item Runs straight from the mains: \hl{cheap and rugged}.
\end{itemize}
\gofurther{the Kloss formula $T/T_\text{max}=2/(s/s_k+s_k/s)$ is an approximation; vector control makes an induction motor behave like a servo motor [Leonhard, 2001].}
\end{column}
\end{columns}
\end{frame}

% ---------------------------------------------------------------------
\begin{frame}{Variable speed with a frequency inverter (V/f)}
\begin{columns}[T]
\begin{column}{.56\textwidth}
\begin{tikzpicture}
\begin{axis}[mkplot,width=7.6cm,height=5.0cm,xmin=0,xmax=1550,ymin=0,ymax=1.15,
  xlabel={speed $n$ (rpm)},ylabel={$T/T_\text{max}$},scaled x ticks=false,
  xticklabel style={/pgf/number format/fixed},legend pos=north west,legend columns=3]
\addplot[mkSky,thick,domain=0:599,samples=150]{2/((600-x)/375+375/(600-x))}; \addlegendentry{\SI{20}{\hertz}}
\addplot[mkMid,thick,domain=0:1049,samples=150]{2/((1050-x)/375+375/(1050-x))}; \addlegendentry{\SI{35}{\hertz}}
\addplot[mkBlue,thick,domain=0:1499,samples=150]{2/((1500-x)/375+375/(1500-x))}; \addlegendentry{\SI{50}{\hertz}}
\addplot[mkOrange,thick,dashed,domain=0:1550,samples=2,forget plot]{0.4};
\end{axis}
\end{tikzpicture}\\
{\scriptsize\color{mkGray} Voltage scaled with frequency: same curve shape, shifted along the speed axis.}
\end{column}
\begin{column}{.42\textwidth}
\small
\begin{itemize}\setlength\itemsep{2pt}
\item A \term{frequency inverter} (VFD) makes 3-phase power at any frequency $f$.
\item $n_s=60f/p$ moves with $f$ $\Rightarrow$ speed control.
\item Keep $V/f$ constant so the field (and max torque) stays the same.
\end{itemize}
\begin{defbox}[Example: conveyor]
\small
4-pole motor, conveyor needs \SI{1000}{rpm}: set $f\approx\frac{1000\cdot2}{60}+$ slip $\approx\SI{35}{\hertz}$.
\end{defbox}
\end{column}
\end{columns}
\end{frame}

% ---------------------------------------------------------------------
\begin{frame}{Which motor where? Choosing for an automation task}
\centering
\small
\renewcommand{\arraystretch}{1.15}
\begin{tabularx}{\textwidth}{@{}l*{3}{>{\centering\arraybackslash}X}@{}}
\toprule
 & \textbf{Stepper} & \textbf{Induction + VFD} & \textbf{PMSM servo}\\
\midrule
Position feedback & none (counts steps) & usually none & \textcolor{mkGreen}{\textbf{encoder}}\\
Accuracy / dynamics & good at low speed & \textcolor{mkRed}{speed only} & \textcolor{mkGreen}{\textbf{excellent}}\\
Torque at high speed & \textcolor{mkRed}{low} & good & good\\
Cost & \textcolor{mkGreen}{\textbf{low}} & \textcolor{mkGreen}{\textbf{low}} & high\\
Typical task & 3D printers, small gantries, lab automation & conveyors, pumps, fans, spindles & robot joints, pick-and-place, packaging\\
\bottomrule
\end{tabularx}

\vspace{4mm}
\begin{columns}[T]
\begin{column}{.5\textwidth}
\begin{keybox}[Decision questions]
\small
Must it stop at exact positions? Must it react fast? Is the load known and constant? What may it cost?
\end{keybox}
\end{column}
\begin{column}{.46\textwidth}
\small
\begin{itemize}\setlength\itemsep{1pt}
\item A production line often uses \hl{all three}.
\item The robot uses servos; the conveyor feeding it an induction motor.
\end{itemize}
\end{column}
\end{columns}
\end{frame}

% ---------------------------------------------------------------------
\begin{frame}{Session 1.4 -- Check your understanding}
\begin{columns}[T]
\begin{column}{.62\textwidth}
\begin{keybox}[Check your understanding]
\begin{enumerate}\setlength\itemsep{6pt}
\item A stepper with 200 steps/rev drives a belt with \SI{40}{\milli\metre} travel per revolution. How far does one full step move the belt?
\item The controller sent 1000 pulses, but the load jammed briefly. Where does the controller think the axis is -- and where is it?
\item A 4-pole induction motor on \SI{50}{\hertz} turns at \SI{1455}{rpm}. What is the slip?
\end{enumerate}
\end{keybox}
\end{column}
\begin{column}{.34\textwidth}
\small
\textbf{Hint:} step angle $=360^\circ/200$; open loop means no measurement; $n_s=60f/p$.
\end{column}
\end{columns}
\end{frame}

% ---------------------------------------------------------------------
\begin{frame}{Session 1.4 -- Summary and key terms}
\begin{columns}[T]
\begin{column}{.55\textwidth}
\small
\begin{itemize}
\item Stepper: moves by pulses, no sensor; torque drops with speed; steps can be lost.
\item Microstepping smooths motion and improves resolution.
\item Induction: rotating field induces rotor currents; needs slip; cheap and rugged.
\item A VFD changes $f$ and so the speed; constant $V/f$ keeps the torque.
\end{itemize}
\begin{keybox}[Remember]
$1.8^\circ$ = 200 steps/rev, \quad $n_s=60f/p$, \quad $s=(n_s-n)/n_s$
\end{keybox}
\end{column}
\begin{column}{.42\textwidth}
\begin{defbox}[Key terms]
\small
\term{Open loop}: no measurement of the result.\\
\term{Microstepping}: sine-shaped staircase currents.\\
\term{Pull-out torque}: max torque before losing steps.\\
\term{Slip}: how much slower the rotor turns than the field.\\
\term{Breakdown torque}: max torque of an induction motor.\\
\term{VFD}: variable-frequency drive (inverter).
\end{defbox}
\end{column}
\end{columns}
\end{frame}

% =====================================================================
\section*{Wrap-up}
% =====================================================================
\begin{frame}{Day summary: W1D1}
\begin{columns}[T]
\begin{column}{.5\textwidth}
\small
\begin{defbox}[Energy conversion and actuators]
Current $+$ field $=$ force; motion $+$ field $=$ voltage; $k_t=k_e$.\\
Electric for joints, pneumatic for grippers, hydraulic for heavy duty.
\end{defbox}
\begin{defbox}[Motor types]
DC (brushes), BLDC (6 steps, Hall), PMSM (sine, encoder), stepper (pulses), induction (slip).
\end{defbox}
\end{column}
\begin{column}{.48\textwidth}
\small
\begin{defbox}[Motor model]
$v=Ri+L\dot i+k_e\omega$, \; $J\dot\omega=k_ti-B\omega-T_L$\\
Fast $\tau_e=L/R$, slow $\tau_m=RJ/(k_tk_e)$.
\end{defbox}
\begin{defbox}[Key numbers]
$P=T\omega$, \; $f_\text{el}=pn/60$, \; $n_s=60f/p$, \; $F=pA$.
\end{defbox}
\end{column}
\end{columns}
\vspace{2mm}
\begin{keybox}[One sentence]
\small
A robot joint is an \hl{energy converter} (motor) whose behaviour two simple ODEs describe -- and whose type we choose by accuracy, dynamics and cost.
\end{keybox}
\end{frame}

% ---------------------------------------------------------------------
\begin{frame}{Preview of W1D2: Actuators II + Introduction to Control}
\begin{columns}[T]
\begin{column}{.55\textwidth}
\small
\renewcommand{\arraystretch}{1.15}
\begin{tabular}{@{}lp{5.3cm}r@{}}
\toprule
\textbf{Code} & \textbf{Session} & \textbf{Time}\\
\midrule
1.5 & Robot joint actuators: gearboxes, reflected inertia, backlash, series-elastic & 1\,h\\
1.6 & Linear actuators and end-effectors & 1\,h\\
1.7 & Motor sizing and selection for a robot joint & 1\,h\\
2.1 & Modelling dynamic systems & 1\,h\\
2.2 & Laplace transform and transfer functions & 1\,h\\
\bottomrule
\end{tabular}
\end{column}
\begin{column}{.42\textwidth}
\begin{notebox}
\small
\textbf{Bring along from today:}\\
$P=T\omega$, $T=k_ti$, the two motor equations, $\tau_e$ and $\tau_m$, and the gravity example. Tomorrow we put a gearbox between motor and link and ask: \hl{which motor is big enough?}
\end{notebox}
\end{column}
\end{columns}
\end{frame}

% ---------------------------------------------------------------------
\begin{frame}{Figure credits}
\scriptsize
\renewcommand{\arraystretch}{1.1}
\begin{tabularx}{\textwidth}{@{}l l X l@{}}
\toprule
\textbf{Slide} & \textbf{File (Wikimedia Commons)} & \textbf{Author} & \textbf{License}\\
\midrule
Robot joint chain & KUKA industrial robotic arm \dots{} The Beach \dots{} 02.jpg & Oleg Yunakov & CC BY-SA 4.0\\
Fluid power & A pneumatic cylinder (Max. press 1.0MPa).jpg & Jonashtand & CC BY-SA 4.0\\
Actuator comparison & Hydraulic cylinder.JPG & S.J. de Waard & CC BY 2.5\\
Brushed DC motor & Motor disassembeled.JPG & Jfer91 & CC BY-SA 3.0\\
BLDC motor & Stator Winding of a BLDC Motor.jpg & Medvedev & CC BY-SA 3.0\\
PMSM & Floppy drive spindle motor open.jpg & Sebastian Koppehel & CC BY 3.0\\
Stepper motor & Positioner Stepper Motor parts.jpg & tony\_duell & CC BY 2.0\\
Induction motor & Wirnik by Zureks.jpg & Zureks & CC BY-SA 3.0\\
\bottomrule
\end{tabularx}
{\tiny Source pages: \texttt{https://commons.wikimedia.org/wiki/File:<file name>} (see \texttt{figures/CREDITS.tsv}).}
\vspace{2mm}

All other drawings, plots and diagrams: own work for this course (TikZ / pgfplots / circuitikz).
\end{frame}

% ---------------------------------------------------------------------
\begin{frame}{References}
\scriptsize
\begin{itemize}\setlength\itemsep{1pt}
\item A.\,E. Fitzgerald, C. Kingsley, S.\,D. Umans, \textit{Electric Machinery}, 6th ed., McGraw-Hill, 2003.
\item A. Hughes, B. Drury, \textit{Electric Motors and Drives}, 4th ed., Newnes, 2013.
\item R. Krishnan, \textit{Electric Motor Drives: Modeling, Analysis, and Control}, Prentice Hall, 2001.
\item R. Krishnan, \textit{Permanent Magnet Synchronous and Brushless DC Motor Drives}, CRC Press, 2010.
\item N. Mohan, \textit{Electric Drives: An Integrative Approach}, MNPERE, 2003.
\item W. Leonhard, \textit{Control of Electrical Drives}, 3rd ed., Springer, 2001.
\item T. Kenjo, A. Sugawara, \textit{Stepping Motors and Their Microprocessor Controls}, 2nd ed., Oxford Univ. Press, 1994.
\item B. Siciliano, L. Sciavicco, L. Villani, G. Oriolo, \textit{Robotics: Modelling, Planning and Control}, Springer, 2009.
\item J.\,J. Craig, \textit{Introduction to Robotics: Mechanics and Control}, 4th ed., Pearson, 2018.
\item W. Bolton, \textit{Mechatronics: Electronic Control Systems in Mechanical and Electrical Engineering}, 7th ed., Pearson, 2019.
\item B. Armstrong-H\'elouvry, P. Dupont, C. Canudas de Wit, ``A survey of models, analysis tools and compensation methods for the control of machines with friction'', \textit{Automatica} 30(7), 1994.
\item IEC 60034-1, \textit{Rotating electrical machines -- Part 1: Rating and performance}; IEC 61800 series, \textit{Adjustable speed electrical power drive systems}.
\end{itemize}
\end{frame}
\end{document}
